\documentclass[11pt,a4paper]{article}
\usepackage[T1]{fontenc}
\usepackage[utf8]{inputenc}
\usepackage{lmodern,amsmath,amssymb,amsthm,mathtools}
\usepackage[margin=25mm,headheight=14pt]{geometry}
\usepackage{microtype,enumitem,booktabs,longtable,array,xcolor,xurl,fancyhdr,needspace}
\definecolor{linkblue}{RGB}{26,64,99}
\usepackage[colorlinks=true,allcolors=linkblue,bookmarksnumbered=true]{hyperref}
\hypersetup{pdftitle={Choice principles after the OpenAI partition-principle model: complete research results},pdfauthor={Frederico Cancado},pdfsubject={Complete preliminary research record, 7 October 2026}}
\pagestyle{fancy}\fancyhf{}
\fancyhead[L]{\small Choice principles: complete research record}
\fancyhead[R]{\small 7 October 2026}
\fancyfoot[C]{\thepage}
\setlength{\parskip}{3pt}
\setlist{itemsep=3pt,topsep=5pt}
\setcounter{tocdepth}{2}
\setlength{\emergencystretch}{2em}
\newcommand{\ZF}{\mathsf{ZF}}
\newcommand{\ZFA}{\mathsf{ZFA}}
\newcommand{\ZFC}{\mathsf{ZFC}}
\newcommand{\AC}{\mathsf{AC}}
\newcommand{\ACWO}{\mathsf{AC}_{\mathrm{WO}}}
\newcommand{\PP}{\mathsf{PP}}
\newcommand{\NDS}{\mathsf{NDS}}
\newcommand{\FB}{\mathsf{FB}}
\newcommand{\SVC}{\mathsf{SVC}}
\newcommand{\Con}{\operatorname{Con}}
\newcommand{\Ord}{\operatorname{Ord}}
\newcommand{\calP}{\mathcal P}
\newcommand{\inj}{\hookrightarrow}
\newcommand{\onto}{\twoheadrightarrow}
\newtheorem{theorem}{Theorem}[section]
\newtheorem{lemma}[theorem]{Lemma}
\newtheorem{proposition}[theorem]{Proposition}
\newtheorem{corollary}[theorem]{Corollary}
\theoremstyle{definition}\newtheorem{definition}[theorem]{Definition}
\theoremstyle{remark}\newtheorem{remark}[theorem]{Remark}
\title{\textbf{Choice principles after the OpenAI\protect\\ partition-principle model}\protect\\[9pt]\large Complete research results and verification record}
\author{Frederico Can\c{c}ado}
\date{7 October 2026\protect\\\small Preliminary research report, version 2026-10-07.1}
\begin{document}
\maketitle
\begin{abstract}
This report consolidates the results obtained in the present investigation of weak choice and cardinal ordering. Starting from OpenAI's symmetric model of the partition principle without choice, it records a finite coinitial cardinal basis theorem, cardinal well-foundedness and finite antichains, a native order-extension obstruction to the Boolean Prime Ideal Theorem, explicit compactness counterexamples, exact cardinal arithmetic and an all-base exponent-collapse phenomenon. It also presents audited parameter variants, Ramsey spectra, and a different extension obtained by forcing a linear order on the native partial order. The latter preserves substantial choice and covering properties, but its partition principle and global cardinal classification remain unresolved. Earlier permutation-model criteria, positive local constructions and counterexamples are retained with their precise scopes. Proofs and proof sketches are paired with a public supporting record. The new arguments have AI-assisted mathematical checks; they have not received independent human peer review. The report separates established source theorems, deductions, preliminary extensions and open questions, without claiming global priority.
\end{abstract}
\paragraph{Attribution.} The construction solving the original PP-versus-AC problem belongs to OpenAI \cite{OpenAI}. The present report concerns additional consequences and model investigations. Substantial OpenAI Codex assistance was used for proof development, checking, literature searches and writing. Human authorship does not imply that independent human verification has already occurred.
\paragraph{Scope of ``complete.''} This is a consolidated account of the mathematical results saved through the reporting cut-off, including negative results and unresolved obligations. It does not assert completion of the research programme or of the paused five-hour investigation. Longer derivations and independent AI audit records are distributed in the accompanying source archive. Earlier exploratory statements are superseded by the statuses in this report. The original imported conversation is preserved privately, unchanged; it is not reissued as a new theorem paper.
\clearpage\tableofcontents\clearpage
\section{Reading the results and keeping the models distinct}
\label{sec:overview}
\subsection{Notation and status conventions}
All unqualified mathematical implications are over classical $\ZF$. We write $X\preccurlyeq Y$ for an injection $X\inj Y$, $X\cong Y$ for a bijection, and $X\prec Y$ for $X\preccurlyeq Y$ with no reverse injection. These comparisons concern cardinalities; they do not require a globally chosen system of cardinal representatives. The Hartogs number $h(X)$ is the least ordinal not injecting into $X$. The Lindenbaum number $\ell(X)$ is the least nonzero ordinal not a surjective image of $X$, with harmless empty-set conventions handled separately. Under $\ACWO$, the largest well-ordered cardinal injecting into an infinite set $X$ exists; it is denoted $w(X)$.

\begin{description}[style=nextline]
\item[Source theorem.] A result in a cited external manuscript or published work. OpenAI's construction remains an explicit dependency.
\item[Audited deduction.] An argument reconstructed or developed in this project and checked mathematically by separate AI agents and/or the coordinating agent. This is not a claim of external peer review or full formalization.
\item[Conditional result.] A proved implication with additional stated hypotheses whose realization or preservation may remain open.
\item[Unresolved here.] No proof or counterexample has been obtained in this investigation. This label does not assert that the world literature treats the problem as open.
\end{description}

The partition principle is
\[
\PP:\quad X\onto Y\ \Longrightarrow\ Y\inj X.
\]
The injection need not be a section of the given surjection. Every surjection having a section is equivalent to $\AC$, and must not be substituted for $\PP$. Choice for all ordinal-indexed families is $\ACWO$; classical implications give $\PP\Rightarrow\ACWO\Rightarrow\mathsf{DC}$. Choice for arbitrary families of finite or countable \emph{members} is a different assertion from choice for finite or countable \emph{index sets}.

$\mathsf{BPI}$ denotes the Boolean Prime Ideal Theorem, equivalently the existence of an ultrafilter on every nontrivial Boolean algebra. $\mathsf{OP}$ says every set can be linearly ordered. $\mathsf{OEP}$ says every partial order has a linear extension. The implication $\mathsf{BPI}\Rightarrow\mathsf{OEP}\Rightarrow\mathsf{OP}$ is classical; the two order principles must not be identified. The ordinary Ramseyan partition principle is denoted $\mathsf{RPP}$ and is defined in the cardinal chapter; it is not the Ramsey property of a group or normal filter.

\subsection{Three settings}
\begin{enumerate}
\item \textbf{The original pure model $W$.} This is the symmetric model in the pinned OpenAI source. The ground $V$ satisfies $\ZFC$. The distinguished non-well-orderable set $A$ consists of Cohen subsets of a regular ordinal $\kappa$. In the original construction $\kappa=\omega_1$ and the ground parameter is $\mu=((2^{\aleph_0})^+)^V$. The monoid $T$ has ground cardinality $\mu$. For this original model, the audited internal identification is $\mu=\omega_2^W$.
\item \textbf{The finite-order extension $N=W[G]$.} Here $G$ is generic for finite linear orders on subsets of $A$ extending its native partial order. The generic order exists in $N$, while it provably does not exist in $W$. Thus the original order-extension obstruction is removed. This does not establish global $\mathsf{BPI}$ or preserve the original image classification.
\item \textbf{Ordinary permutation models with atoms.} Earlier work used a faithful group $G$ and an effective normal filter of subgroups to build $\ZFA$ models. Its group criteria and counterexamples concern that setting. A pure symmetric extension is not an ordinary atom permutation model; the relevant collapse and transfer arguments cannot be silently moved between them.
\end{enumerate}

The source revision throughout is
\begin{center}\small\path{adc7f1241b42e322a6451854ab7e4b4c146bf78a}.\end{center}
The source supplies $\SVC(A)$, $\ACWO$, $A$ non-well-orderable, and the property
\[
Q(A):\quad\text{every non-well-orderable surjective image of $A$ is bijective with $A$.}
\]
The anchors are Proposition 3.6, Proposition 3.7, Proposition 3.8, Theorem 6.3 and Lemma 7.1 of \cite{OpenAI}. Section 8 supplies the ordinary first-order relative-consistency bridge. None of our applications replaces this bridge by an assumption that a transitive model or an inaccessible cardinal exists.

\subsection{Main outcomes at a glance}
\begin{longtable}{>{\raggedright\arraybackslash}p{.20\linewidth}>{\raggedright\arraybackslash}p{.72\linewidth}}
\toprule
Setting & Recorded outcome\\\midrule\endhead
Original $W$ & $\PP+\ACWO+\mathsf{DC}+\FB+\NDS+\neg\AC$; all set-sized cardinal antichains finite, but their finite sizes unbounded.\\
Original $W$ & $\mathsf{KWP}_1$, $\mathsf{OP}$, choice for finite members, and choice for arbitrary families of countable members; $\neg\mathsf{DC}_{\omega_1}$.\\
Original $W$ & Native partial order has no linear extension; consequently $\neg\mathsf{OEP}$ and $\neg\mathsf{BPI}$. Explicit Boolean and graph witnesses have size $A$.\\
Original $W$ & Exact cardinal normal forms and coordinate comparison; set-sized suprema and nonempty infima; all infinite $X$ satisfy $X^3\cong X^2$.\\
Original $W$ & For every non-well-orderable $X$ and every set $Y$, $Y^X\cong Y^{w(X)}$. In particular $\mu\prec A$ while $Y^\mu\cong Y^A$ for every $Y$.\\
Variants & Regular-parameter adaptations separate ordinal-comparability and Ramsey thresholds while retaining the strong original package and failure of BPI.\\
Extension $N$ & $\ACWO+\mathsf{DC}+\neg\AC+\neg\mathsf{DC}_{\omega_1}$; old cardinals and cofinalities preserved; $\SVC(A)$ and $\SVC(\mathbb R)$; a non-well-orderable continuum.\\
Extension $N$ & $A\inj\mathbb R$, $A\onto\mathbb R$, and $A^\omega\cong\mathbb R$; the two displayed directions do not give $\mathbb R\inj A$. Every generic open interval and final segment has size $A$.\\
Unresolved & General $\PP\Rightarrow\NDS$; a different model of $\PP+\mathsf{BPI}+\neg\AC$; $\PP,Q(A),\FB,\NDS,\mathsf{BPI}$ in $N$; and $\mathbb R^N\inj A$.\\
\bottomrule
\end{longtable}

\subsection{Historical motivation}
The partition principle goes back to the early development of choice principles; a historical account and modern formulation appear in \cite{daSilva,BM}. Howard and Tachtsis studied no decreasing sequence of cardinals \cite{HT}, and Tachtsis studied its position in the weak-choice hierarchy \cite{Tachtsis2024}. Blass and Kulshreshtha distinguish 24 well-foundedness formulations and identify them under $\PP$ \cite{BK}; that equivalence does not prove any of them from $\PP$ alone.

Feldman--Orhon, with Blass's appendix, show that any fixed finite bound on cardinal antichain size forces choice \cite{FO}. The assertion that each antichain is finite has no fixed bound and is different. Karagila explicitly records the corresponding question \cite{Karagila}; Tachtsis studies related infinite-family comparability \cite{Tachtsis2018}. Our use of the label $\FB$ is descriptive. The finite-product order argument is classical well-quasi-order theory \cite{Higman,Rado}; the application to the cited construction is the additional deduction.

\section{Finite bases, NDS and cardinal antichains}\label{defs}

All arguments take place in classical $\ZF$. Write $X\preccurlyeq Y$ if there is an injection $X\inj Y$, and $X\prec Y$ if $X\preccurlyeq Y$ and $Y\not\preccurlyeq X$. Write $X\cong Y$ for the existence of a bijection. The disjoint union is denoted by $\sqcup$.

\begin{definition}
The following are statements about sets.
\begin{enumerate}[label=(\roman*)]
\item $\NDS$: there is no sequence $(X_n)_{n<\omega}$ with $X_{n+1}\prec X_n$ for every $n$.
\item The finite-antichain principle: every set family $\mathcal F$ whose distinct members have incomparable cardinalities is finite.
\item $\FB$: for every nonempty set family $\mathcal F$, there is a nonempty finite $\mathcal F_0\subseteq\mathcal F$ such that
\begin{equation}\label{FBdef}
 \forall X\in\mathcal F\ \exists Y\in\mathcal F_0\quad Y\preccurlyeq X.
\end{equation}
\end{enumerate}
No simultaneous choice of all the injections in~\eqref{FBdef} is part of the assertion. Likewise, $\NDS$ quantifies over sequences of sets, not sequences already equipped with injection witnesses. This is the formulation $\mathsf{WF}_{i3}^{i}$ of \cite{BK}; distinct formulations must not be interchanged in bare $\ZF$.
\end{definition}

For statements phrased directly about cardinalities, one may use Scott representatives: the set of all representatives of the least possible rank. If $\mathcal C$ is a set of Scott cardinalities, its union is a set of representatives. Applying $\FB$ to that union avoids choosing one set for each cardinality.

\begin{definition}\label{package}
Fix a nonempty set $A$. The \emph{homogeneous-image hypotheses} are:
\begin{enumerate}[label=(\Alph*)]
\item\label{S} $\SVC(A)$: every nonempty set $X$ is the image of a surjection $A\times\eta\onto X$ for some ordinal $\eta$.
\item\label{W} $\ACWO$: every ordinal-indexed family of nonempty sets has a choice function.
\item\label{Q} Every non-well-orderable surjective image of $A$ is bijective with $A$.
\end{enumerate}
Non-well-orderability of $A$ itself is required for the consistency application, but not for the abstract finite-basis implication.
\end{definition}

\begin{theorem}[Central implication]\label{central}
In $\ZF$, the hypotheses of Definition~\ref{package} imply $\FB$. Hence every sequence of sets has a forward-comparable pair, $\NDS$ holds, and every set-sized cardinal antichain is finite.
\end{theorem}

The proof is given in Subsections~\ref{normal}--\ref{finitebasis}. The resulting application is
\begin{equation}\label{newcons}
 \Con(\ZF)\ \Longrightarrow\
 \Con(\ZF+\PP+\ACWO+\FB+\NDS+\neg\AC),
\end{equation}
with finiteness of all set-sized cardinal antichains in the same model. The construction in \cite{OpenAI} is the external input for~\eqref{newcons}.

\subsection{Extracting the cardinal normal form}\label{normal}

\begin{lemma}\label{NF}
Under Definition~\ref{package}, every set $X$ has a representation
\begin{equation}\label{normalform}
 X\cong(A\times\alpha)\sqcup\beta,
 \qquad \alpha,\beta\in\Ord.
\end{equation}
The two parameters may be chosen to be initial ordinals; their particular values need not be unique.
\end{lemma}
\begin{proof}
The empty set uses $\alpha=\beta=0$. Otherwise fix one surjection $s:A\times\eta\onto X$. Put
\[
 B_i=s[A\times\{i\}],\qquad
 T_i=B_i\setminus\bigcup_{j<i}B_j\quad(i<\eta).
\]
The $T_i$ partition $X$. Each nonempty $T_i$ is a surjective image of $A$: fix one $t\in T_i$, keep $s(a,i)$ when it belongs to $T_i$, and send all other arguments to $t$. This proves existence for one slice at a time; it does not choose default elements simultaneously.

Let $I$ consist of those $i<\eta$ for which $T_i$ is not well-orderable. For $i\in I$, hypothesis~\ref{Q} supplies a bijection $A\to T_i$; for $i\notin I$, take a well-order of $T_i$, including the empty relation when the slice is empty. The candidate bijections lie in $\calP(A\times X)$ and the candidate orders in $\calP(X\times X)$. Thus $\ACWO$ selects witnesses from a set-sized, well-orderably indexed family.

The union of the non-well-orderable slices is bijective with $A\times I$, hence with $A\times\alpha$ for an ordinal $\alpha$. The remaining slices have a well-orderable union: order by slice index first, then by the selected well-order inside each slice. That union is bijective with an ordinal $\beta$. The two pieces are disjoint and cover $X$. Replacing each ordinal by its initial ordinal preserves the representation.
\end{proof}

The slicing mechanism is an extraction of the method already used in \cite[Lemma 7.1]{OpenAI}, which in turn credits Ryan-Smith's local-to-global argument \cite{RyanSmith}. The next finite-basis step is the additional application needed here.

For each $X$, choose its parameters \emph{definably}: let $\alpha(X)$ be the least ordinal admitting any representation~\eqref{normalform}, and let $\beta(X)$ be the least ordinal witnessing it for that $\alpha(X)$. These exist because any one witness bounds a nonempty set in which to minimize. Replacement collects the pairs for any set family. No class choice or simultaneous selection of normal-form bijections is being made. Arbitrary representations need not be unique; only the chosen least pair is canonical.

\subsection{A finite basis and its consequences}\label{finitebasis}

\begin{lemma}[Ordinal pairs, in $\ZF$]\label{ordpairs}
Every nonempty set $R\subseteq\Ord^2$ has a nonempty finite subset $R_0$ such that every pair in $R$ is coordinatewise above a pair in $R_0$.
\end{lemma}
\begin{proof}
Choose the lexicographically least residual pair, delete its coordinatewise upper cone, and repeat while a pair remains. The rule is unique, so recursion on $\omega$ requires no choice. If it never stopped, the selected first coordinates would be nondecreasing: removal cannot introduce a smaller first coordinate. Since every later pair survived deletion of each earlier upper cone, its second coordinate would be strictly smaller than the preceding second coordinate. This would be an infinite descending sequence of ordinals. The process therefore stops after finitely many selections. Every deleted pair lies above a selected pair, proving the assertion.
\end{proof}

\begin{proof}[Proof of the finite-basis assertion in Theorem~\ref{central}]
For a nonempty set family $\mathcal F$, apply Lemma~\ref{ordpairs} to
\[
 R=\{(\alpha(X),\beta(X)):X\in\mathcal F\}.
\]
Choose one member of $\mathcal F$ for each of the finitely many basis pairs. This is finite choice, a theorem of $\ZF$. Let $\mathcal F_0$ be the resulting finite family. If the profile of $Y\in\mathcal F_0$ is coordinatewise below that of $X\in\mathcal F$, the coordinate inclusions give
\[
 (A\times\alpha(Y))\sqcup\beta(Y)
 \inj (A\times\alpha(X))\sqcup\beta(X).
\]
Composing with two normal-form bijections gives $Y\preccurlyeq X$. This proves~\eqref{FBdef}. The injections need only exist separately.
\end{proof}

\begin{corollary}\label{seq}
Under $\FB$, every sequence $(X_n)_{n<\omega}$ has indices $i<j$ such that $X_i\preccurlyeq X_j$. In particular, $\NDS$ holds.
\end{corollary}
\begin{proof}
Apply $\FB$ to the range of the sequence. Each member of its finite basis has a least occurrence index. Choose $j$ larger than all these indices. A basis member injects into $X_j$, and its least index $i$ is smaller than $j$.

If the sequence were strictly decreasing, finitely many downward injections would give $X_j\preccurlyeq X_{i+1}$. Composing contradicts $X_i\not\preccurlyeq X_{i+1}$. Only finitely many witnesses are used, not countably many chosen injections.
\end{proof}

\begin{corollary}\label{antichain}
Under $\FB$, every set-sized cardinal antichain is finite. More generally, every nonempty set family represents finitely many minimal cardinalities, and each of its cardinalities lies above one of them.
\end{corollary}
\begin{proof}
For an antichain, a basis member can inject into another member only when they represent the same cardinality. Hence the finite basis accounts for every cardinality in the antichain. For a general family, remove redundant cardinalities from a finite basis, leaving its minimal ones. They remain coinitial. If a family cardinality lay strictly below one of them, another basis member would lie below that smaller cardinality, contradicting minimality in the finite basis. Conversely every minimal family cardinality must coincide with a basis cardinality.
\end{proof}

There can still be infinitely many distinct sets representing a single minimal cardinality. Corollary~\ref{antichain} concerns the number of cardinalities, not the number of such representatives. The Scott-set formulation follows by applying $\FB$ to the union of the sets of representatives.


\paragraph{Relationship with well-foundedness.} Over $\ZF$, $\FB$ is equivalent to well-foundedness of the injection order on each set of cardinalities together with finiteness of all cardinal antichains. For the converse, the minimal members of a nonempty family form a finite antichain, and well-foundedness supplies a minimal member below each given member. Under $\mathsf{DC}$, absence of an infinite strict descending sequence is equivalent to this set-sized well-foundedness: a nonempty downward set with no minimal member would yield a descending sequence by dependent choice. Consequently, over $\ZF+\mathsf{DC}$, $\FB$ is equivalent to $\NDS$ plus finite cardinal antichains. The qualification by $\mathsf{DC}$ cannot simply be erased.

\paragraph{Relative consistency.} Applying the theorem to the stronger source package and retaining the source's non-well-orderable $A$ yields $\Con(\ZF)\Rightarrow\Con(\ZF+\PP+\FB+\NDS+\neg\AC)$. Under the same $\PP$, the well-foundedness equivalence of \cite[Proposition 9.8]{BK} gives the other 23 formulations. These are consequences in this model; neither $\PP\Rightarrow\NDS$ nor $\NDS\Rightarrow\PP$ is established generally here.

\clearpage% Retrospective section. No preamble; requires amsmath, amssymb and url/hyperref.
% Mathematical records below concern original W, not its ordinary order extension.
\section{Cardinal order, exponentiation, and choice in the original model}
\label{sec:complete-cardinals}

This section collects the results established for the \emph{original}
symmetric model $W$, and, where explicitly indicated, its separately audited
parameter variants. It does not transfer these conclusions to the later
ordinary forcing extension. The foundational construction, its Partition
Principle and support lemmas are due to OpenAI~\cite{OpenAI}; the additional
calculations below are the research deductions reported here. They have
received separate mathematical checks by AI agents, not external human peer
review or a complete Lean formalization. Historical qualifications are given
at the end of the section.

\subsection{The generator and the exact cardinal records}

Write $X\leq Y$ for an injection and $X\cong Y$ for a bijection. In the
source's general ground $V\models\mathsf{ZFC}$, put
\[
 \kappa=\omega_1^V,\qquad \mu=((2^{\aleph_0})^+)^V.
\]
The set $A\subseteq\mathcal P(\kappa)$ consists of the distinct Cohen rows.
The source proves $\mathsf{SVC}(A)$, choice for ordinal-indexed families
$\mathsf{AC}_{\mathrm{WO}}$, and the classification
\[
 \text{every non-well-orderable image of $A$ is bijective with $A$}. \tag{Q}
\]
Together these give $\mathsf{PP}$ and $\neg\mathsf{AC}$ in $W$.
The same source locality implies that a function on an ordinal with values
in a ground set belongs to the Cohen intermediate extension. Thus ground
cardinals at least $\mu$ remain cardinals in $W$.

\paragraph{The action and absorption.}
The generic diagrams define an internal free right action of the ground
monoid $T$ on $A$. For each $a$, the map $t\mapsto a\cdot t$ is injective;
its image $aT$ is a \emph{chart}, of cardinal $\mu$. Every ordinal-indexed
subset of $A$ lies in one chart. The latter assertion follows by copying a
putative outside value freshly while fixing the supporting diagram: two
compatible conditions would assign different labels to the same ordinal
index. Hence the largest well-orderable cardinal injecting into $A$ is
\[
                         w(A)=\mu,\qquad h(A)=\mu^+.
\]
Here $h$ is the Hartogs number. No global inverse-label map $A\to I$ is
asserted to belong to $W$.

Each map $a\mapsto a\cdot t$ is onto. Define a predecessor diagram by
$q(tx)=p(x)$ and put its remaining coordinates freshly; exact amalgamation
makes such predecessors dense. For $t\ne1$, every fiber over an anchor
$a_0$ is non-well-orderable. To see the key step, normalize an alleged
ordinal enumeration to a condition $(d,p)$ containing its support and the
Cohen support, with $a_0=A_{p(ce_0)}$, where $e_0\ne1$ is the distinguished
monoid element. Take $q_{e_0}=p$. The source probe clause at
$(e_0c,t,\{1\})$ supplies $n$ with
\[
                   nt=e_0ce_0,\qquad e_0\not\sqsubseteq n,
\]
where $u\sqsubseteq v$ means $uz=v$ for some $z$.
Thus $A_{q(n)}$ is in the fiber but outside the root diagram. Decide its
enumeration index and copy both the stronger diagram and its Cohen
condition freshly over the root. Exact diagram overlap and disjoint
Cohen supports off the root give the contradiction.

Fibers for distinct $t$ are disjoint by freeness in the monoid coordinate.
Each nonempty fiber is an image of $A$ by a retraction, so (Q) identifies
it with $A$. Choosing these bijections over the well-ordered set
$T\setminus\{1\}$ gives $A\times\mu\leq A$. Consequently
\[
                  A\times\mu\cong A,\qquad A^2\cong A.       \tag{1}
\]
For the second equality, first use the raw normal form below for $A^2$.
Its well-orderable subsets have well-orderable projections of size at most
$\mu$, so both ordinal parameters are at most $\mu$ and are absorbed by
$A\times\mu$. This avoids assuming square absorption in advance.

\paragraph{Normal form and order.}
An $\mathsf{SVC}(A)$ presentation $A\times\eta\twoheadrightarrow X$ splits
$X$ by its least presenting ordinal coordinate. By (Q), each resulting
non-well-orderable piece is a copy of $A$; the other pieces are
well-orderable. Ordinal-indexed choice assembles the bijections, yielding
$X\cong(A\times\alpha)\sqcup\beta$.
Using (1), the exact records are
\[
 \begin{array}{ll}
 (0,\lambda):&\text{a well-orderable cardinal }\lambda,\\
 (a,e),\ \mu\leq a\leq e:
   & A\times a\quad\text{if }a=e,\\[-2pt]
   &(A\times a)\sqcup e\quad\text{if }a<e.
 \end{array}                                                    \tag{2}
\]
All nonzero coordinates here are ordinal cardinals. For the second kind,
$w(X)=e$, $h(X)=e^+$, and
\[
             A\leq X,\qquad e\leq X\leq A\times e.
\]
Under $\mathsf{AC}_{\mathrm{WO}}$, the Lindenbaum number agrees with the
Hartogs number. The least non-well-orderable cardinal is therefore $|A|$.

The records determine the \emph{entire injection order}:
\[
            (a,e)\leq(c,f)\quad\Longleftrightarrow\quad
                         a\leq c\ \hbox{and}\ e\leq f,        \tag{3}
\]
including the well-orderable records with first coordinate zero.
For necessity in the first coordinate, an injection between non-WO
normal forms cannot send an entire source $A$-slice into the target's
ordinal tail. Ordinal-indexed choice picks one point from each source
slice whose image lies in the target's $A$-part. These are distinct,
so the source coefficient is at most the largest WO cardinal of that
part, namely the target coefficient. Necessity for $w$ is immediate.
Conversely embed the $A$-part and the ordinal part using the two coordinate
inequalities, then combine the embeddings using doubling of the target.
The WO cases reduce to the definition of $w$.

\subsection{Lattice structure, finite bases, and finite powers}

Equation (3) gives a precise cardinal lattice, with necessary size
qualifications: \emph{every set-sized family} of cardinalities has a
supremum, and every \emph{nonempty set-sized family} has an infimum.
They are obtained by coordinatewise suprema and minima in (2).
The resulting pairs are again realizable records. Binary distributive
laws follow from those for minimum and maximum in two ordinal chains.
There is no largest cardinal and no empty infimum; this is not an
unqualified claim of a complete lattice or of complete distributivity.
Scott cardinals can be used throughout to avoid choosing representatives.

The finite-basis theorem proved elsewhere in this report says that every
nonempty family has a finite coinitial family under injections. It yields
$\mathsf{NDS}$ and the absence of infinite antichains of cardinalities.
It does \emph{not} bound all antichains by one fixed natural number. For
any positive $n$, the $n$ records
\[
 \bigl(\mu^{+i},\mu^{+(2n-i)}\bigr),\qquad i<n,
\]
form an antichain: the first coordinates increase and the second decrease.
In bare ZF the exact order-theoretic equivalence is
$\mathsf{FB}\leftrightarrow\mathsf{WF}_{\mathrm{card}}+\mathsf{FAC}$,
where $\mathsf{WF}_{\mathrm{card}}$ is minimal-element well-foundedness
on sets of Scott cardinals and $\mathsf{FAC}$ means every such antichain
is finite. With DC it becomes
$\mathsf{FB}\leftrightarrow\mathsf{NDS}+\mathsf{FAC}$.
The representative and sequence distinctions should not be suppressed
when DC is absent. Moreover, the already proved PP and NDS give all
24 cardinal well-foundedness variants in Blass--Kulshreshtha's
Definition~4.1, using their equivalence under PP
\cite[Proposition 9.8]{BK}. This equivalence does not prove NDS from
PP alone.

For a fixed ordinal cardinal $\lambda$, the assertion
$\mathsf W_\lambda$ that every cardinal is comparable with $\lambda$
has the exact threshold
\[
                W\models\mathsf W_\lambda
                       \quad\Longleftrightarrow\quad\lambda\leq\mu.
                                                                  \tag{4}
\]
Every non-WO set contains $A$, proving the positive direction; $A$ is
incomparable with every ordinal cardinal above $\mu$, proving the negative.
The usual surjective version agrees under PP. In particular full
comparability still fails, as it must when AC fails.

Every infinite set absorbs a countable multiplier and doubles. For
non-WO $X$ with record $(a,e)$, direct multiplication gives
\[
                 X^n\cong A\times e\quad(2\leq n<\omega).     \tag{5}
\]
Indeed $A\times e\leq X^2\leq A^2\times e^2$, and multiplying the
normal form by $A\times e$ leaves the same cardinal. Thus
$X^3\cong X^2$ holds for every infinite $X$, but $X^2\cong X$ does not:
$X=A\sqcup\mu^+$ has record $(\mu,\mu^+)$ and its square has record
$(\mu^+,\mu^+)$. The question about universal square/cube stability is
historical; Rubin--Rubin's 1985 introduction already discusses it. The
bounded search does not establish its intervening publication status.

There is also a finite bound, depending on $X$, on iterated power-set
roots: some $N<\omega$ satisfies
\[
                    \mathcal P^n(Y)\cong X\quad\Longrightarrow\quad n\leq N.
                                                                  \tag{6}
\]
Represent every possible root by a subset of $X$, using iterated singleton
maps. The resulting set of roots has a finite coinitial basis
$Y_1,\ldots,Y_k$, with unique depths $n_1,\ldots,n_k$; uniqueness follows
from Cantor's theorem. If a root $Y$ had depth $m>\max n_i$, some
$Y_i\leq Y$ would give
$X\cong\mathcal P^{n_i}(Y_i)\leq\mathcal P^{n_i}(Y)
 <\mathcal P^m(Y)\cong X$, impossible. In particular every cardinal is
a finite power-set iterate of one which is not a power-set cardinal.
There is no bound uniform in $X$.

\subsection{Supported names and exact exponentiation}

For ground ordinals $b,\lambda$, write
\[
             E_V(b,\lambda)=|\{f\in V:f:\lambda\to b\}|^V.
\]
This is a ground cardinal; unmarked exponentials are computed in $W$.
The decisive local calculation, for $b\geq2$ and ordinal cardinals
$\lambda\geq\mu$, is
\[
 b^A\cong A\times E_V(b,\mu),\qquad
 b^\lambda\cong A\times E_V(b,\lambda).                         \tag{7}
\]
The coefficients are at least $(2^\mu)^V$ and are preserved cardinals.

Here are the ingredients ensuring that (7) counts \emph{all internal}
well-orderable families. Fix a supporting diagram $p$. Normalize names
to globally total $A$-to-$b$ functions, using the zero function whenever
the original value has the wrong type. The canonical graph name uses
only label/ordinal pair names and the ordinary forcing relation. It is
equivariant, keeps the support, and has a fixed rank bound. Thus these
names form a ground set, rather than a class of unbounded-rank names.
For each $t\in T$ choose a ground diagram $q^t$ with $(q^t)_t=p$.
At the record $(t,s,d)$, with $d$ a Cohen condition on $D_{q^t}$, record
$\xi<b$ exactly when
\[
                 (d,q^t)\Vdash\tau(A_{q^t(s)})=\xi.
\]
There are $\mu$ positions, and at most one value per position. The code
is a \emph{partial value function}, so there are at most
$(|b|^V+1)^\mu=E_V(b,\mu)$ codes in the ground.

Equal codes force equality below $(1,p)$. Otherwise put an input where
the values differ into a larger diagram $r$, write $p=r_t$, restrict
the Cohen witness to $D_r$ by source locality, and transport it to $q^t$
by a permutation fixing $D_p$. It produces different values at the same
record, a contradiction. Ground choice now selects one representative
per code. Their ordinal enumeration is HS with support $D_p$.
Any internal WO family of functions has one supported enumeration;
normalizing its coordinates with the same support embeds the whole family
into that representative list by least indices. Ground functions on
one chart give the matching lower bound. For ordinal domain $\lambda$
the records are $(d,\xi)$, $\xi<\lambda$, and both local-decision clauses
give the analogous count. There is no maximum-principle assumption,
choice indexed by $A$, or automorphism of the actual pure model in this
argument. Finally both function sets are non-WO and square-idempotent;
(2)--(5) turn their exact WO bounds into (7).

\paragraph{Universal exponent collapse.}
For every non-WO $X$, set $e=w(X)$. Then for \emph{every set} $Y$,
\[
                         Y^X\cong Y^e.                         \tag{8}
\]
More explicitly,
\[
 \begin{aligned}
 \mathcal P(X)&\cong A\times E_V(2,e)\cong\mathcal P(e),\\
 b^X&\cong A\times E_V(b,e) &&(b\geq2\text{ ordinal}),\\
 Y^X&\cong A\times E_V(f,e) &&(Y\text{ non-WO},\ f=w(Y)).
 \end{aligned}                                                   \tag{9}
\]
The same formulas hold with domain $e$.

For the proof, two-point coding and $A\leq\mathcal P(\kappa)$ give
$A^e\cong\mathcal P(e)$. The normal form also gives
$\kappa\times X\cong X$, hence $A^X\cong\mathcal P(X)$.
Sandwich an ordinal-base power between $b^e$ and $b^{A\times e}$.
With $\delta=E_V(b,\mu)$, (7) calculates the upper endpoint as
\[
 (b^A)^e\cong(A\times\delta)^e
 \cong A^e\times\delta^e\cong A\times E_V(b,e),
\]
since ground arithmetic has $E_V(\delta,e)=E_V(b,e)$.
The lower endpoint has the same cardinal. Its case $b=2$ proves the
power-set formula before that formula is used for non-WO bases.
For non-WO $Y$, use $A\leq Y$, $f\leq Y\leq A\times f$ for the upper
bound. To inject $A\times(f^e)^V$ into $Y^e$, reserve one coordinate for
the $A$ point and $e$ further coordinates for the ground function,
using $e+1\cong e$. This provides the matching lower bound in both
$Y^e$ and $Y^X$. Empty and singleton bases give the trivial equalities
because $X$ and $e$ are nonempty. This exhausts all internal bases.

Consequently
\[
                  \mu<A,\qquad \forall Y\ (Y^\mu\cong Y^A).    \tag{10}
\]
The statement asserts pointwise existence of bijections. It asserts
neither a natural transformation between function-space functors nor
one class-sized choice of all bijections.

\subsection{Complete arithmetic table over the constructible ground}

For the separately audited $L$-ground variants $M_{m,n}$, put
$\kappa=\aleph_m$, $\mu=\aleph_n$, where $1\leq m<n<\omega$.
All alephs are preserved. The original construction over $L$ is $M_{1,2}$.
The following tables concern these original symmetric models.
The records of power sets are
\[
\begin{array}{c|c}
 |X|&|\mathcal P(X)|\text{ as a record}\\ \hline
 \text{finite }r&(0,2^r)\\
 \text{infinite WO }\lambda<\kappa&(0,\lambda^+)\\
 \text{WO }\kappa\leq\lambda<\mu&(\mu,\mu)\\
 \text{WO }\lambda\geq\mu&(\lambda^+,\lambda^+)\\
 \text{non-WO }(a,e)&(e^+,e^+)
\end{array}                                                       \tag{11}
\]
For the small ordinal plateau, the local Cohen forcing has size $\mu$
and the $\kappa^+$ chain condition. Nice names for subsets of $\lambda$
number at most $\mu^{\max(\kappa,\lambda)}$, equal under ground GCH to
$\mu$ for $\kappa\leq\lambda<\mu$ and to $\lambda^+$ for
$\lambda\geq\mu$. One supported enumeration bounds every WO family;
a chart and the ground power set supply the lower bounds. Below $\kappa$
there are no new subsets of ordinals. In particular, in $M_{1,2}$,
\[
 \mathbb R\cong\aleph_1,\qquad
 \mathcal P(\aleph_1)\cong A,\qquad
 \mathcal P(A)\cong\mathcal P(\aleph_2)\cong A\times\aleph_3.
\]
This is a ground-GCH calculation, not GCH inside the choiceless model.

For completeness write $E(b,\lambda)=E_L(b,\lambda)$. Apart from empty
and singleton cases, the \emph{entire exponent table} is
\[
\begin{array}{c|c|c}
 |Y|&|X|&|Y^X|\text{ as a record}\\ \hline
 \text{WO }b\geq2&\text{WO }2\leq\lambda<\kappa&(0,E(b,\lambda))\\
 \text{WO }b\geq2&\text{WO }\lambda\geq\kappa&(\theta,\theta)\\
 \text{non-WO }(a,f)&\text{WO }\lambda\geq2&(E(f,\lambda),E(f,\lambda))\\
 \text{WO }b\geq2&\text{non-WO }(c,e)&(E(b,e),E(b,e))\\
 \text{non-WO }(a,f)&\text{non-WO }(c,e)&(E(f,e),E(f,e))
\end{array}                                                       \tag{12}
\]
where $\theta=\max(\mu,E(b,\lambda))$.
For $X=\varnothing$ the result is $1$, including $0^0=1$;
for $X\ne\varnothing$ and $Y=\varnothing$ it is $0$;
a singleton base always gives $1$; a singleton exponent returns $Y$
with its original record. This last exception matters when $a<f$.

To justify the rows not already covered by (9), ordinal-valued functions
of length $\lambda<\kappa$ are ground functions. Also
$A^\lambda\cong A$ for $0<\lambda<\kappa$: a sequence has range in one
chart and its $T$-coordinate sequence is ground, giving
$A\times(T^\lambda)^V\twoheadrightarrow A^\lambda$; now
$\mu^{<\kappa}=\mu$, (1), and PP apply. For $\lambda\geq\kappa$,
nice function names number at most
$(\max(\mu,b)^\lambda)^L=\max(\mu,E(b,\lambda))$.
This gives the WO-base row as above. The non-WO-base upper bound is
$A^\lambda\times f^\lambda$. A matching lower bound uses two coordinates
if $\lambda$ is finite, since $f^\lambda=f$ in the ground, and the
reserved-coordinate coding if $\lambda$ is infinite.

The ground coefficients in (12) are explicit. For positive finite
$\lambda$, $E(b,\lambda)$ is the usual finite power when $b$ is finite
and is $b$ when $b$ is infinite. For infinite $\lambda$ and $b\geq2$,
\[
 E(b,\lambda)=
 \begin{cases}
  \lambda^+,&b\leq\lambda,\\
  b,&\lambda<b\text{ and }\lambda<\operatorname{cf}^L(b),\\
  b^+,&\operatorname{cf}^L(b)\leq\lambda<b.
 \end{cases}                                                     \tag{13}
\]
The final case concerns singular $b$; it follows from K\"onig's theorem
and $b^\lambda\leq2^b=b^+$.

\subsection{Choice consequences and their exact boundaries}

In addition to $\mathsf{PP}$, $\mathsf{AC}_{\mathrm{WO}}$,
$\mathsf{FB}$ and $\mathsf{NDS}$, the original model satisfies ordinary
DC, $\mathsf{KWP}_1$, the Ordering Principle, and choice for arbitrary
families of nonempty finite sets. The classical implication
$\mathsf{AC}_{\mathrm{WO}}\Rightarrow\mathsf{DC}$ is Jensen's theorem
as presented in Jech~\cite[Theorem 8.2]{Jech}; an earlier apparent source
conflict was resolved by reading the actual PDF. For $\mathsf{KWP}_1$,
PP turns a presentation into $X\leq A\times\eta$, which embeds into a
power set of an ordinal by tagged coding. Lexicographic order on that
power set linearly orders $X$. Least elements of finite subsets then
give finite-member choice. This does not assert that every prescribed
partial order has a linear extension: the separate obstruction in this
report proves failure of that stronger principle and of BPI.

The source itself explicitly states the dual Cantor--Bernstein principle
$\mathsf{CSB}^*$ and Weak Partition Principle as consequences of PP.
With $X\leq^*Y$ meaning $X=\varnothing$ or $Y\twoheadrightarrow X$,
these say respectively
\[
 X\leq^*Y\ \&\ Y\leq^*X\Rightarrow X\cong Y,
 \qquad X\leq^*Y\Rightarrow\neg(Y<X).
\]
They are not newly extracted results of this report.

\paragraph{Choice for countable members.}
A further saved proof, now independently audited twice, gives
\[
 \mathsf C(\infty,\leq\aleph_0):\quad
 \text{every family of nonempty countable sets has a choice function}.
                                                                  \tag{14}
\]
This is Howard--Rubin Form 85A. The family of indices is arbitrary;
its \emph{members} are countable, so (14) differs from countable choice.
For one nonempty countable $S\subseteq A\subseteq\mathcal P(\omega_1)$,
its pairwise first-difference ordinals are countably many and therefore
bounded. Let $\alpha(S)<\omega_1$ be the least ordinal whose restriction
map separates $S$. Source preservation gives
$\mathcal P(\alpha)^W=\mathcal P(\alpha)^V$ for every $\alpha<\omega_1$.
One fixed ground family of well-orders of these power sets selects the
least restriction and hence a unique member of $S$. No simultaneous
choice of enumerations of $S$ was used. Given arbitrary $X$, choose one
embedding $X\leq A\times\eta$ and first take the least ordinal slice of
a countable subset, then this selector on $A$. For an arbitrary family,
apply the construction just once to its union. The proof is for original
$W$; new reals prevent an automatic transfer to the later extension.

\paragraph{Ordinary versus higher dependent choice.}
Despite DC, the original model and the audited parameter family fail
$\mathsf{DC}_{\omega_1}$. Order $A$ natively by
$a\leq_Tb$ iff $a\cdot t=b$ for some $t\in T$. Every ordinal-indexed
subset has a common lower bound by chart locality, and every point has
a strict predecessor by surjectivity of the $e_0$-translation.
There is no strictly decreasing $\omega_1$-sequence: localize its whole
range to one chart, lift to $T$, and apply the source's strictly
increasing real-valued height. This would give a decreasing
$\omega_1$-sequence of reals, impossible because its disjoint successive
intervals inject into the rationals. Thus the tree of descending
sequences of countable ordinal length is countably closed and has no
terminal node but has no $\omega_1$-branch. Equivalently, the relation
``strictly below every member of the preceding countable set'' directly
violates the conclusion of Howard--Rubin's Form 44.
Increasing only the Cohen parameter does not remove this height obstruction.

\subsection{Ramsey principles: ordinary truth and stronger separations}

Here RPP means the \emph{Ramseyan Partition Principle}: every two-coloring
of pairs of an infinite set has an infinite homogeneous subset. This is
not PP and not the group-filter Ramsey condition used in FM analyses of
BPI. Ordinary RPP holds in $W$: DC gives a countably infinite subset,
and the ordinary finite-arity Ramsey theorem on $\omega$ is provable in
ZF. More generally, $\mathsf{NDS}$ implies that every infinite set is
Dedekind-infinite: otherwise finite injective lists from an infinite
Dedekind-finite set produce a strictly descending sequence by excluding
successively shorter lists. Thus both the PP/DC and NDS routes give RPP.
These implications are classical weak-choice context~\cite{Blass1977,BK}.

The stronger cardinal demands differ. Let $\mathrm{NR}_\lambda$ say
that every two-coloring of pairs of every non-WO set has a homogeneous
subset of cardinality $\lambda$. Throughout the audited parameter family,
\[
                 \mathrm{NR}_\lambda
                          \quad\Longleftrightarrow\quad\lambda\leq\kappa
                 \quad(\lambda\text{ infinite WO}).            \tag{15}
\]
For the positive part pull a coloring back to a $\kappa^+$-subset of
$A$. An inner $L[C]$ containing its ordinal code and the required small
sequence enumerations satisfies the appropriate Erd\H{o}s--Rado theorem;
its witness remains in $W$. The first-difference parity coloring on
$A\subseteq\mathcal P(\kappa)$ gives the negative part: fresh-copy
arguments make every homogeneous set WO, and a Cohen-condition
$\Delta$-system argument bounds its cardinal by $\kappa$. Each color
also has a homogeneous set of size $\kappa$. The arbitrary-ground pair
proof uses the separately audited collapse of $2^{<\kappa}$ to $\kappa$
in $W$; the $L$-ground cases need only their unchanged smaller power sets.

The color-count boundary is also exact. For every ordinal cardinal
$1\leq\eta<\kappa$, every $\eta$-coloring of pairs of a non-WO set has a
$\kappa$-sized homogeneous subset. At $\eta=\kappa$ the coloring
$\{a,b\}\mapsto\min(a\mathbin\triangle b)$ has no homogeneous triple:
three distinct binary bits cannot be pairwise different at one coordinate.
Thus $\kappa$ is the least WO number of colors permitting such a failure.

For the $L$-ground family $M_{m,n}$, let
$\mathrm{NR}^{r,\mathrm{fin}}_\lambda$ assert that every coloring of
$[X]^r$ into any positive finite number of colors, for every non-WO $X$,
has a homogeneous set of cardinality $\lambda$. Its exact spectrum is
\[
 M_{m,n}\models\mathrm{NR}^{r,\mathrm{fin}}_\lambda
 \quad\Longleftrightarrow\quad
 \lambda\leq\aleph_{\max(m-r+2,0)},\qquad r\geq2.             \tag{16}
\]
Here $\lambda$ ranges over infinite WO cardinals. Inner-model Erd\H{o}s--Rado
supplies the positive bound; a generalized Sierpi\'nski coloring and
first-difference stepping up give the negative bound~\cite{ErdosRado,BHT}.
The proof checks that homogeneous sets in the stepping-up construction
are well-ordered by lexicographic order or its reverse, rather than
assuming this in advance. The triple obstruction uses \emph{four colors}.
The exact two-color triple question remains unresolved here; the audit
of simple binary coarsenings does not settle it.

For example, $M_{1,3}$ and $M_{2,3}$ have the same
$w(A)=\aleph_3$, $h(A)=\aleph_4$, and the same entire threshold (4).
Their pair homogeneous thresholds are $\aleph_1$ and $\aleph_2$;
their finite-color triple thresholds are $\aleph_0$ and $\aleph_1$.
Both retain PP, FB, NDS, $\mathsf{AC}_{\mathrm{WO}}$, DC and failure of
$\mathsf{DC}_{\omega_1}$ and AC. This separates the stronger Ramsey
statements within that shared package. It is not a claim that the
classical colorings or partition calculus are new.

\subsection{Comparability, cancellation, and historical boundaries}

All power sets in $W$ have comparable cardinalities. An infinite set
doubles, so its power set squares to itself. Every non-WO power set
therefore has record $(\theta,\theta)$, and these form a chain. A WO
power set must have a WO base of cardinal below $\kappa$; otherwise it
contains $A\leq\mathcal P(\kappa)$. Such smaller powers have ground
enumerations below $\mu$, so lie below every non-WO power. Thus
\[
 \mathsf{PC}:\qquad
  \forall X,Y\ [\mathcal P(X)\leq\mathcal P(Y)
                  \ \text{or}\ \mathcal P(Y)\leq\mathcal P(X)]. \tag{17}
\]
Bare $\mathsf{PC}+\neg\mathsf{AC}$ is prior art: Glazer's publicly dated
30 September 2025 answer proves it in Cohen's first model~\cite{Glazer}.
The joint PP/FB/DC package and exact coefficients are the application
under discussion here; no exhaustive novelty guarantee is claimed.
In ZF, PC also implies $\mathsf{KWP}_1$: compare $\mathcal P(X)$
with $\mathcal P(\theta)$ for $\theta=h(\mathcal P(X))$. The latter
cannot inject into the former, so $X\leq\mathcal P(X)\leq\mathcal P(\theta)$.
This observation appears in Wojowu's comment on the cited question.

The Injective Continuum Function assertion
$\mathsf{ICF}:\mathcal P(X)\cong\mathcal P(Y)\Rightarrow X\cong Y$
fails already in ZFC at a permissible continuum plateau, and is not a
consequence of AC. Here it fails more specifically at the strict pair
$\mu<A$, by (10). The order-reflecting variant
$\mathcal P(X)\leq\mathcal P(Y)\Rightarrow X\leq Y$ fails as well.
For finite Kinna--Wagner principles,
\[
          \mathsf{ICF}+\mathsf{KWP}_{n+1}
                           \Rightarrow\mathsf{KWP}_{n}\quad(n<\omega),
                                                                  \tag{18}
\]
so ICF plus any finite $\mathsf{KWP}_n$ implies AC.
Apply $\mathsf{KWP}_{n+1}$ to $\mathcal P(X)$ to get
$\mathcal P(X)\leq\mathcal P(B)$ with $B=\mathcal P^n(\theta)$,
$\theta$ infinite. Since $B$ doubles,
$\mathcal P(X\sqcup B)\cong\mathcal P(B)$; ICF gives $X\leq B$.
Here $\mathsf{KWP}_0=\mathsf{AC}$. No transfinite version or priority
claim is asserted~\cite{KS,ICFquestion}. The corresponding finite
surjective principles $\mathsf{KWP}^{*}_n$ imply $\mathsf{KWP}_{n+1}$
by taking inverse images under the presentation surjection. Thus ICF
combined with any finite $\mathsf{KWP}^{*}_n$ implies AC as well.

For clarity, distinguish two all-base assertions: (EC) if
$Y^X\cong Y^Z$ for every $Y$, then $X\cong Z$; and (SES) if $X<Z$,
some $Y$ satisfies $Y^X<Y^Z$. These are descriptive labels, not asserted
standard names. Both hold in ZFC. For infinite cardinals $\alpha<\beta$,
choose $\nu=\beth_{\alpha^+}$, a strong limit of cofinality $\alpha^+$;
then $\nu^\alpha=\nu<\nu^\beta$ by bounded ranges and K\"onig's theorem.
Finite cases use base $2$. Equation (10) refutes both in $W$.

Bare comparable all-base plateaux have an older mechanism. Mutual
surjections induce injections between function spaces by precomposition
at every base. From an infinite Dedekind-finite set $D$, let $S$ be its
finite injective lists. Deleting the first entry, with the empty list
sent to an extra point, maps $S$ onto $S\sqcup1$; a retraction maps back.
But $S<S\sqcup1$, so their powers agree at every base. Hamkins's 2012
answer records this list construction; the all-base conclusion is the
immediate precomposition extension~\cite{Hamkins}. Thus bare failure of
EC or SES is not a new consistency phenomenon. The distinction in (10)
is its compatibility with PP, DC and dual CSB, and its one well-orderable
exponent; it cannot arise from unequal mutually surjective sets.
For the logical record, ZF proves
$\mathrm{EC}\Rightarrow\mathrm{SES}$: on a comparable nonempty pair,
function extension gives the forward injections for every nonempty base;
if none were strict, EC would give equality of the exponents. An empty
smaller exponent is separated by base $2$. It also proves
$\mathrm{EC}\Rightarrow\mathsf{CSB}^*$, by the precomposition argument,
and either EC or SES implies that every infinite set is
Dedekind-infinite, by the finite-list counterexample just given.
No converse or general equivalence with AC is established here.
Finally strict \emph{reflection},
$p^x<p^y$ and $p\ne0\Rightarrow x<y$, is a different statement, listed
as equivalent to AC by Lindenbaum--Tarski~\cite[No.~82(A5)]{LT}.
Its direction and quantifiers do not identify the strength of EC or SES.

The general implication $\mathsf{PP}\Rightarrow\mathsf{NDS}$ has not been
proved by these model calculations. Nor do they prove PP alone implies
FB, the exact arithmetic, or the stronger Ramsey statements. Earlier
phase notes that left BPI in this concrete model undecided are superseded
by the separate negative result in this report. Positive BPI in a different
PP model, and the global choice classification of the later extension,
remain distinct questions.

\paragraph{Proof-file map and verification scope.}
Detailed records are in \nolinkurl{evidence/original_model/}:
\nolinkurl{T_ACTION_FIBERS_AND_SQUARE_COEFFICIENT_DRAFT.txt};
\nolinkurl{T_ACTION_INDEPENDENT_AUDIT_AND_DC_OMEGA1.txt};
\nolinkurl{EXACT_EXPONENTIALS_AND_UNIVERSAL_EXPONENT_COLLAPSE_AUDIT.txt};
\nolinkurl{COMPLETE_GCH_GROUND_EXPONENT_TABLE.txt};
\nolinkurl{FB_NDS_FAC_EXACT_CHOICE_AUDIT.txt};
\nolinkurl{COUNTABLE_MEMBER_CHOICE_SECOND_INDEPENDENT_AUDIT.txt};
\nolinkurl{GLOBAL_NON_WELL_ORDERABLE_RAMSEY_THRESHOLD.txt};
\nolinkurl{HIGHER_ARITY_ROOT_INDEPENDENT_AUDIT.txt}; and
\nolinkurl{ALL_BASE_EXPONENT_CANCELLATION_LITERATURE_AUDIT.txt}.
The lattice and root-depth proofs are in
\nolinkurl{evidence/finite_basis/PROOFS_cardinal_lattice_and_power_roots.txt}.
These files preserve source dependencies and search limits. The new results
in this section are mathematical deductions, not a certification of
novelty or a claim that all of them have been formally verified.

\clearpage% Consolidation only. No preamble. Requires amsmath, amssymb and hyperref.
% Saved audits, not new research; 7 October 2026.
\section{Boolean prime ideals, compactness, and model variations}
\label{sec:bpi-variants}

This section records consequences of the released OpenAI construction
and adaptations checked in the present research. The construction itself
is credited to OpenAI, at commit
\href{https://github.com/openai/math/commit/adc7f1241b42e322a6451854ab7e4b4c146bf78a}{\texttt{adc7f124}}.
The new arguments below have mathematical drafts and separate internal
AI-agent audits. That is distinct from external peer review and from a
formal proof of their full statements. Their novelty has not been
certified. Older working notes saying that BPI in this model is
``undecided'' are superseded by the six-cycle theorem below.

Write $X\simeq Y$ for equipotence, and $X\preccurlyeq Y$ for the existence
of an injection. The relevant established package in the original model
$W$ is
\[
 \mathsf{ZF}+\mathsf{PP}+\mathsf{AC}_{\rm WO}+\mathsf{DC}
 +\mathsf{FB}+\mathsf{NDS}+\neg\mathsf{AC}.
\]
Here $\mathsf{FB}$ is the finite cardinal basis property defined earlier,
not an abbreviation for finite choice. The additional structural inputs
are $\mathsf{SVC}(A)$, the quotient dichotomy $Q(A)$, and the least
non-well-orderable cardinal $|A|$. Every set embeds into $A\times\eta$
for an ordinal $\eta$, every non-well-orderable set contains a copy of
$A$, and $A^{<\omega}\simeq A^2\simeq A$. The right action of the ground
monoid $T$ on $A$ is free on each orbit. The native order is
\[
 a\leq_T b\quad\Longleftrightarrow\quad
 (\exists t\in T)\ a\cdot t=b.
\]
Left cancellation and conicality of $T$ make this a partial order.
Its graph is an internal hereditarily symmetric object; its definition
does not require treating the generic filter as an internal parameter.

\subsection{A partial order with no linear extension}

\textbf{Audited theorem.}
The native partial order on $A$ has no linear extension in $W$.
Consequently
\[
 W\models\neg\mathsf{OEP}+\neg\mathsf{BPI},
\]
where OEP says that every partial order has a linear extension and BPI
says that every nontrivial Boolean algebra has a prime ideal. The same
argument applies to the parameter variants described below. The
Ordering Principle still holds: every set can be linearly ordered,
but an arbitrary prescribed partial order need not be extended.

\emph{Proof sketch with the compatibility step.}
The source monoid has a distinguished $e\ne1$ and $\mu$ many probes
$m_\xi$ satisfying
\[
 m_\xi e=e,\qquad m_\xi\ne1,\qquad
 m_\xi T\cap m_\zeta T=eT\quad(\xi\ne\zeta).
\]
It also realizes every valid finite prescription of principal-cone
overlaps, with exactly the prescribed identifications. A diagram is a
ground injection $q:T\to I$; write $q_t(x)=q(tx)$ and $D_q=q[T]$.
Cohen conditions mention fewer than $\mu$ coordinate labels.

Suppose a name for a linear extension $L$ is supported by $D_p$ and
$(c_0,p)$ forces its asserted properties. Restrict $c_0$ to $D_p$.
Choose $q$ with $q_e=p$ and two probes $a,b$. The source local-decision
lemma lets $q$ decide their comparison after the Cohen generic is fixed.
A Cohen strengthening $d$, restricted to $D_q$, therefore forces, after
possibly exchanging $a,b$,
\[
 q\Vdash A_{q(a)}<_L A_{q(b)}.
\]
This fixes the diagram before choosing the following native step.

The cones $bm_\xi T\setminus eT$ are pairwise disjoint. Hence some
probe $t\ne1$ has
\[
 \operatorname{supp}_I(d)\cap q[btT\setminus eT]=\varnothing.
\]
Set $h=bt$. Then $he=e$, $aT\cap hT=eT$, and $b<_T h$.
On three copies, indexed modulo $3$, prescribe
\[
                   (i,hx)\sim(i+1,ax)\qquad(x\in T).
\]
This prescription is valid. Indeed, a two-step passage through the
middle copy must lie in $aT\cap hT=eT$; if $ax=hy=ez$, cancellation
and $ae=he=e$ give $x=y=ez$. Thus the indicated relation is already
transitive, its common root has identical parameters, and no copy
collapses onto itself.

Place these copies as diagrams $q_i$, all with $(q_i)_e=p$, and transport
$d$ to conditions $d_i$ by coordinate permutations fixing $D_p$.
The $d_i$ agree on the root. Off the root, each outgoing overlap is
silent in the corresponding $d_i$ by the choice of $t$. Their union
is therefore a valid Cohen condition. This is the essential step:
a bare swap of two diagrams would not guarantee compatible Cohen bits.
The diagram merger gives a common extension of the $q_i$. With
$x_i=A_{q_i(a)}$ and $y_i=A_{q_i(b)}$, it forces
\[
 x_i<_L y_i<_T x_{i+1}\quad(i\bmod3),
\]
and hence a strict six-cycle in $L$, a contradiction.

Forcing automorphisms here transport names and conditions. They are
not asserted to be automorphisms of the actual pure ZF universe.
The implication $\mathsf{BPI}\Rightarrow\mathsf{OEP}$ is classical
propositional compactness; see Jech, \emph{The Axiom of Choice} (1973),
Example 2.3.2, p.~19. Jech's \S7.2 already uses finite symmetry cycles
in models with OP and without OEP. The model-specific contribution here
is the silent probe cone and its exact Cohen-compatible merger.

\emph{Verification scope.}
The auxiliary \path{SixCycleAudit.lean} proves validity of this cyclic
overlap using the pinned source's actual
\texttt{OverlapPrescription.Valid}, exact root intersections, absence
of within-copy collapse, and compatibility of copied partial assignments.
Its compilation reported no \texttt{sorryAx}. It does \emph{not}
formalize the full symmetric-forcing argument or the conclusion
$W\models\neg\mathsf{BPI}$.
The principal records are
\path{NATIVE_ORDER_EXTENSION_SIX_CYCLE_PROOF.txt},
\path{INDEPENDENT_AUDIT_SIX_CYCLE_OEP_FAILURE.txt}, and
\path{SIX_CYCLE_LEAN_SCOPE.txt}.

\subsection{Explicit Boolean-algebra and graph witnesses}

\textbf{Boolean algebra.}
There is a nontrivial Boolean algebra $B_{\rm ord}$ of cardinality
exactly $|A|$ with no prime ideal. Use variables $p_{ab}$, $a,b\in A$,
and impose
\[
 \neg p_{aa},\quad p_{ab}\vee p_{ba}\ (a\ne b),\quad
 \neg(p_{ab}\wedge p_{ba})\ (a\ne b),\quad
 (p_{ab}\wedge p_{bc})\Rightarrow p_{ac},\quad
 p_{ab}\ (a<_T b).
\]
Take the free Boolean algebra modulo the ideal generated by the
negations of these constraints. Every finite subtheory is satisfiable:
linearly extend the native partial order on its finitely many mentioned
points. Thus the quotient is nontrivial. A prime ideal would give a
homomorphism to $2$, hence a forbidden native linear extension.

Finite terms have cardinality $A$, so $A$ surjects onto $B_{\rm ord}$.
Every well-orderable nontrivial Boolean algebra has a prime ideal in ZF:
recursively include the next enumerated element in a growing proper
ideal whenever this preserves properness. If it cannot be included,
its complement already belongs to the ideal. Limit unions stay proper.
Consequently $B_{\rm ord}$ is not well-orderable, and $Q(A)$ gives
$B_{\rm ord}\simeq A$. It has the least possible injection-cardinality
of any prime-ideal-free nontrivial Boolean algebra. In particular, for
every set $X$ in this model,
\[
 \begin{split}
 X\text{ is well-orderable}\quad\Longleftrightarrow\quad&
 \text{every nontrivial Boolean algebra }B\preccurlyeq X\\[-2pt]
 &\text{has a prime ideal}.
 \end{split}
\]

\textbf{Finite-cone realization and an unavoidable new subset.}
Put $S=[A]^{<\omega}$, $C_F=\{E\in S:F\subseteq E\}$, and
\[
 \mathcal I=\{D\subseteq S:(\exists F\in S)\ D\cap C_F=\varnothing\}.
\]
Canonically linearly extend each finite native order on $E$ by repeatedly
taking its lexicographically least remaining minimal point; call the
result $L_E$. The old sets
\[
 D_{ab}=\{E\in S:a,b\in E\text{ and }a<_{L_E}b\}
\]
satisfy the order constraints modulo $\mathcal I$. The subalgebra
generated by their classes inside $\mathcal P(S)/\mathcal I$ is another
size-$A$ algebra with no prime ideal. The entire quotient need not
have size $A$.

Even an ultrafilter on the \emph{old} algebra $\mathcal P(S)^W$ containing
all $C_F$ defines a native linear extension by
$a<_U b\Longleftrightarrow D_{ab}\in U$. Thus forcing with the positive
part of $\mathcal P(S)^W/\mathcal I$ necessarily adds a new subset of
$A^2$, and therefore of $A$. Moreover, the $A^2$-indexed dense open
sets deciding the Boolean values $[D_{ab}]$ have empty intersection.
In the Boolean completion they witness failure of binary distributivity
indexed by $A$. No claim about countable distributivity follows.
Genericity decides old subsets; it does not automatically give an
ultrafilter on the enlarged final powerset.

\textbf{Graph compactness.}
There is a graph $G$ with $|G|=|A|$ such that every finite subgraph is
$3$-colorable but $\chi(G)=4$. The order theory above is a finitely
satisfiable $3$-CNF theory after repeating literals. Encode Boolean
variables by opposite vertices adjacent to a fixed base vertex in a
triangle of colors $T,F,B$. A clause $x\vee y\vee z$ is equivalent to
\[
 \exists w\;\bigl(\operatorname{NAE}(x,y,w)
 \wedge\operatorname{NAE}(\neg w,z,F)\bigr),
\]
where NAE means ``not all equal.'' Each NAE constraint is represented
by a fresh triangle whose three vertices are respectively adjacent
to the three Boolean terminals. It has an extending $3$-coloring
exactly when NAE holds. A global $3$-coloring would solve the impossible
order theory. Every finite subgraph is colorable by finite
satisfiability, while a canonical $4$-coloring assigns Boolean values
uniformly and colors each gadget greedily. Tagged finite powers of
$A$ give the exact vertex cardinality.

Well-orderable graphs satisfy finite-color compactness in ZF by
transfinite recursion, choosing at each vertex the first color
preserving finite satisfiability. Thus $A$ is again a least possible
counterexample cardinal. In contrast, $2$-color compactness holds in
$W$: finite bipartiteness gives a canonical pair of parity classes in
each nontrivial connected component, and $\mathsf{AC}_2$ selects one
class per component. Here $\mathsf{AC}_2$ follows from OP.
For each fixed finite $k\geq3$, joining this graph with a complete graph on $k-3$ fresh vertices gives a graph whose finite subgraphs are $k$-colorable while its chromatic number is $k+1$. Thus finite-$k$ compactness fails for every such $k$.
The general compactness connection is classical; see de Bruijn--Erd\H{o}s
(1951) and Cowen--Hechler--Mihok (2002).
The explicit gadgets and infinite argument are checked in
\path{EXPLICIT_GRAPH_COMPACTNESS_INDEPENDENT_AUDIT.txt}; the Boolean
proof is in
\path{ORDER_OBSTRUCTION_BOOLEAN_ALGEBRA_AND_FINE_FORCING.txt}.

\subsection{BPI, small violations of choice, and limits of compactness}

Under $\mathsf{SVC}(A)$, BPI is equivalent to the existence of a fine
ultrafilter on $[A]^{<\omega}$, equivalently a prime ideal in
$\mathcal P([A]^{<\omega})/\mathcal I$. This is the Pincus--Blass local
reflection recorded in Ryan-Smith, \emph{Local reflections of choice},
Proposition 3.13. The preceding theorem settles this particular local
test negatively in $W$; it does not prove the abstract implication
$\mathsf{SVC}(A)+\mathsf{AC}_{\rm WO}+Q(A)+\mathsf{BPI}\Rightarrow
\mathsf{AC}$.

The model gives relative nonimplications
\[
 \mathsf{PP}\not\Rightarrow\mathsf{BPI},\qquad
 \mathsf{NDS}\not\Rightarrow\mathsf{BPI},\qquad
 \mathsf{FB}\not\Rightarrow\mathsf{BPI},
\]
and likewise with OEP as conclusion. In the reverse direction, the
classical basic Cohen model satisfies BPI and SVC and has an infinite
Dedekind-finite set. It fails PP, NDS, FB, and Ramsey's principle on
arbitrary infinite sets. This is an old model, not a new construction;
the SVC generator need not be its Cohen set. Tachtsis's 2024 Theorem
6.4 supplies the stronger opposite-direction example
$\mathsf{BPI}+\mathsf{AC}_\omega+\neg\mathsf{NDS}$, hence also
$\neg\mathsf{FB}$ and ordinary infinite homogeneity. Its exact
attribution was verified through Blass--Kulshreshtha's primary
manuscript, Theorem 10.5; the full Tachtsis article was not obtained.
These statements do not decide whether PP implies NDS in arbitrary
ZF models, or whether PP and BPI can coexist with $\neg\mathsf{AC}$.

Several useful subsidiary facts survived the unsuccessful positive-BPI
routes. FB gives every proper filter a member of minimum injection
cardinality; restricting an ultrafilter to that member makes it
cardinal-uniform. This is not a cofinal filter base or an ultrafilter
existence theorem. The native lower cones generate a proper filter
closed under all internal well-ordered intersections and avoiding every
well-orderable subset of $A$. An ultrafilter extending it need not retain
that completeness. In the original model every nonprincipal ultrafilter
on $A\subseteq\mathcal P(\omega_1)$ is countably incomplete: first
push forward by the first disagreement with its coordinatewise
ultrafilter decisions, obtaining a nonprincipal image on $\omega_1$;
then use the ground injection $\omega_1\to\mathcal P(\omega)$ and
repeat, obtaining a nonprincipal image on $\omega$.

Likewise, BPI and PP give a valid small-carrier Henkin lemma:
a finitely satisfiable finitary theory whose symbols inject into $A$
has a model of cardinality at most $A$, provided $A^{<\omega}\simeq A$.
This does not transport the original monoid action through an arbitrary
bijection of carriers. Finite-factor ultraproduct arguments also do
not force strict external cardinal descent: in ZFC every unbounded
finite product $\prod_n k(n)/U$, for nonprincipal $U$ on $\omega$,
has size $2^{\aleph_0}$, even after finitely many iterated logarithms
of the factors. These boundaries are recorded in
\path{STRUCTURAL_BPI_PHASE1.txt},
\path{ANCESTOR_FILTER_COMPACTNESS_AND_ULTRAFILTER_INCOMPLETENESS.txt},
\path{HENKIN_SMALL_CARRIER_AND_STRUCTURE_OBSTRUCTION.txt}, and
\path{ULTRAPRODUCT_AND_QUOTIENT_OBSTRUCTIONS.txt}.

\subsection{Parameter changes and preservation}

\textbf{Audited parameter theorem.}
Let the ground satisfy ZFC. Suppose $\kappa$ is regular uncountable,
$\mu$ is regular, and
\[
 2^{<\kappa}<\mu,\qquad
 \lambda^{<\kappa}<\mu\quad
 \text{for every infinite ground cardinal }\lambda<\mu.
\]
Use the source monoid of size $\mu$, $|I|=\mu^+$, Cohen conditions of
size $<\kappa$, and diagram supports of size at most $\mu$. The audited
adaptation yields the same SVC, ordinal-choice and quotient package,
hence PP, FB, NDS and $\neg\mathsf{AC}$. The quotient proof requires
an explicit repair: bounded arrays indexed below $\omega_1$ become
arrays of size $<\kappa$, controlled by the strengthened no-new-ground-
sequences lemma. The countable germ construction remains countable.
The native six-cycle and real-height obstruction remain valid, giving
$\neg\mathsf{BPI}$ and $\neg\mathsf{DC}_{\omega_1}$ in every variant,
although ordinary DC holds.

Two internal thresholds are distinct. If $W_\lambda$ says that every
set is injection-comparable with the well-ordered cardinal $\lambda$,
then
\[
 w(A)=\mu,\quad h(A)=\mu^+,\qquad
 W_\lambda\ \Longleftrightarrow\ \lambda\leq\mu,
 \qquad
 \mathcal P(\alpha)\text{ is well-orderable}\
 \Longleftrightarrow\alpha<\kappa.
\]
Here $w(A)$ denotes the largest well-ordered cardinal injecting into
$A$. For $\alpha<\kappa$, the entire powerset equals its ground
counterpart. Under $V=L$, write $M_{mn}$ for $\kappa=\aleph_m$,
$\mu=\aleph_n$, $1\leq m<n<\omega$. All ground alephs are preserved:
\begin{center}
\begin{tabular}{c|cc|cc}
Model & $\kappa$ & $\mu$ & $W_{\aleph_3}$ &
$\mathcal P(\aleph_1)$ well-orderable\\\hline
$M_{12}$ & $\aleph_1$ & $\aleph_2$ & no & no\\
$M_{13}$ & $\aleph_1$ & $\aleph_3$ & yes & no\\
$M_{23}$ & $\aleph_2$ & $\aleph_3$ & yes & yes
\end{tabular}
\end{center}
These are specific separations over a common strong package, not
unrestricted independent assignments of two sentences. The consistency
interpretation uses the source's formula-by-formula forcing method over
$L$; ordinary $\operatorname{Con}(\mathsf{ZF})$ suffices, with no
inaccessible cardinal. The complete adaptation is not Lean-formalized.
See \path{LARGER_MU_SECOND_INDEPENDENT_AUDIT.txt}.

Without GCH, put $\theta=(2^{<\kappa})^V$. The ground Cohen forcing is
$\theta^+$-cc and preserves $\kappa$ and its regularity. Every old
binary string of length $<\kappa$ occurs as a block of any one generic
row. The row's block map therefore collapses the old set of such strings,
of size $\theta$, to cardinality $\kappa$ internally in $W$. Thus
\[
 (\kappa^+)^W=(\theta^+)^V,\qquad
 (\kappa^{<\kappa})^W=\kappa.
\]
Ordinal-function locality transfers the relevant preservation from
the Cohen extension, not from the full outer product extension.
The exact no-GCH pair-coloring argument respects this identification;
see \path{NO_GCH_HOMOGENEOUS_BOUND_RECHECK.txt}.

\textbf{Preservation theorem.}
If $W\subseteq N$ have the same ordinals, $N$ adds no subsets of $A$,
and $N$ satisfies $\mathsf{SVC}(A)$, then $N$ retains
$\mathsf{AC}_{\rm WO}+Q(A)+\neg\mathsf{AC}$ and its PP/FB consequences.
Indeed, $A^2\simeq A$ makes every new map on $A$ have an old kernel.
An old non-well-orderable quotient cannot become well-orderable without
well-ordering $A$. Old ordinal-choice then supplies representatives for
kernels of new ordinal-valued maps. SVC reduces general ordinal-choice
to these maps. Ordinary set forcing always preserves SVC by evaluating
a ground presentation of the set of subnames.

This criterion has concrete instances. If the outer ZFC forcing has
size $\chi$, take regular $\lambda>\chi$ and the ground forcing
$\operatorname{Add}(\lambda,\delta)^V$. Easton's small-forcing argument
makes it $<\lambda$-distributive in the outer extension. The audited
intersection lemma $W[K]\cap V[F][H]=W$ then gives no new subsets of
$A^2$ when $\lambda$ exceeds its outer size. Taking $\lambda$ also above
the outer size of $\mathcal P([A]^{<\omega})^W$ preserves the BPI truth
value; here that means preserving its failure. For the original
$L$-ground model one may use $\lambda=\aleph_5^V$ and
$\delta\geq\lambda^+$ with $\delta^\lambda=\delta$. The largest
well-ordered cardinal injecting into $\mathcal P(\lambda)$ becomes
$\delta$, while that powerset stays non-well-orderable. These are
arithmetic variations, not further weak-choice separations by themselves.
See \path{HIGH_FORCING_PRESERVATION_INDEPENDENT_AUDIT.txt}.

\subsection{Exact pair and finite-arity Ramsey spectra}

For an infinite well-ordered cardinal $\lambda$, use the temporary
notation $\mathrm{NR}^{r,\mathrm{fin}}_\lambda$ for: every finite
coloring of $[X]^r$, for every non-well-orderable $X$, has a homogeneous
subset of cardinality $\lambda$. The requirement concerns all sets and
colorings, not only a distinguished witness. Under AC it is vacuous.
Ordinary Ramsey homogeneity with an unspecified infinite size already
follows from DC here.

\textbf{Pairs, without GCH.}
Throughout the audited family, the largest $\lambda$ for binary pair
colorings is exactly $\kappa$. More strongly, every pair coloring with
an ordinal color cardinal $1\leq\eta<\kappa$ on a non-well-orderable
set has a homogeneous subset of size $\kappa$. At $\eta=\kappa$,
the coloring
\[
 c(\{a,b\})=\min(a\mathbin{\triangle}b),\qquad a,b\in A,
\]
has no homogeneous triple: three binary bits cannot be pairwise
different at the same first-difference coordinate. Thus $\kappa$ is
the exact least well-ordered color cardinal permitting a failure of
infinite pair homogeneity on a non-well-orderable set.

For the positive direction, encode a pulled-back coloring on
$\kappa^+$, the enumeration of $\kappa^{<\kappa}$, and ordinal-sized
injection witnesses into a set $C$ of ordinals. The inner model $L[C]$
computes the required cardinals and $2^{<\kappa}=\kappa$ correctly.
Apply the Erd\H{o}s--Rado theorem there, in the form
$(2^{<\kappa})^+\to(\kappa+1)^2_\eta$ for $\eta<\kappa$.
The resulting homogeneous set belongs to $W$. Use
$\kappa^+\preccurlyeq A$ and then the copy of $A$ in every
non-well-orderable set. The precise primary source is
Baumgartner--Hajnal--Todorcevic, \emph{Extensions of the Erd\H{o}s--Rado
Theorem}, Theorem 2.1, p.~4.

For the negative direction, color $\{a,b\}$ by the parity of their
first difference. Fresh-copy locality first makes every homogeneous
subset well-orderable. A Cohen $\Delta$-system argument excludes size
$\kappa^+$: thin conditions to equal row patterns and then prescribe
a fresh first difference of the opposite parity. Both colors have
homogeneous subsets of size $\kappa$, supplied by rows whose first $1$
occurs at increasing coordinates of a fixed parity. In the no-GCH case
the argument uses the preserved ground $\theta^+$, which is the
internal $\kappa^+$. Consequently
\[
 \mathrm{NR}^{2,\mathrm{fin}}_\lambda
 \quad\Longleftrightarrow\quad \lambda\leq\kappa.
\]
Sources and complete proofs are
\path{EXACT_HOMOGENEOUS_CARDINAL_BOUND.txt} and
\path{GLOBAL_NON_WELL_ORDERABLE_RAMSEY_THRESHOLD.txt}.

\textbf{Higher arities over $L$.}
In $M_{mn}$ the exact threshold for all finite color counts is
\[
 \boxed{\quad
 \max\{\lambda:\mathrm{NR}^{r,\mathrm{fin}}_\lambda\}
       =\aleph_{\max(m-r+2,0)},\qquad 2\leq r<\omega.
 \quad}
\]
\begin{center}
\begin{tabular}{c|cccc}
Model & pairs & triples & $4$-tuples & $5$-tuples\\\hline
$M_{12},M_{13}$ & $\aleph_1$ & $\aleph_0$ & $\aleph_0$ & $\aleph_0$\\
$M_{23}$ & $\aleph_2$ & $\aleph_1$ & $\aleph_0$ & $\aleph_0$\\
$M_{34}$ & $\aleph_3$ & $\aleph_2$ & $\aleph_1$ & $\aleph_0$
\end{tabular}
\end{center}
The positive proof applies finite-arity Erd\H{o}s--Rado inside $L[C]$;
all powersets needed for its finite GCH calculation lie strictly below
$\kappa$ and are unchanged. The negative proof begins with Sierpi\'nski's
comparison of a ground ordinal enumeration and binary lexicographic
order, then steps up using consecutive first differences. A tuple is
colored by the preceding color together with an increasing/decreasing
tag, or by an additional nonmonotone color. Finite Ramsey excludes an
infinite homogeneous set of the latter color. DC makes a homogeneous
set in either remaining orientation well-ordered by lex or reverse lex;
its first differences yield a forbidden homogeneous set for the prior
coloring. This avoids an auxiliary well-order of $\mathcal P(\kappa)$.
The negative triple witness uses four colors; successive constructions
use at most $5\cdot2^{r-3}-1$ colors, not claimed optimal.
See \path{HIGHER_ARITY_ROOT_INDEPENDENT_AUDIT.txt} and Erd\H{o}s--Rado
(1956), Theorem 39, pp.~467--469.

\textbf{Binary triples remain a separate question.}
The saved bounds are
\[
\begin{array}{c|ccc}
 &\text{proved positive}&\text{unresolved here}&\text{proved negative}\\\hline
 M_{12},M_{13}&\omega&\omega_1&\omega_2\\
 M_{23}&\omega_1&\omega_2&\omega_3
\end{array}
\]
for homogeneous cardinalities of arbitrary binary triple colorings on
non-well-orderable domains. The binary orientation coloring of adjacent
first differences has exact largest homogeneous cardinality $\kappa$,
in both colors; this gives the stated negative upper bounds, not a
universal positive theorem at $\kappa$. Every one of the sixteen maps
merging the particular four-color triple witness into two colors has
an $\omega_1$-sized homogeneous set in $M_{12}$ and $M_{13}$. Projections
use row patterns; the XOR case uses a countable binary tree of
coordinates labeled by nested real intervals. Thus simple color merging
cannot settle the missing negative result. The classical binary
stepping-up proof uses an auxiliary well-order unavailable internally.
See \path{BINARY_TRIPLE_ROOT_INDEPENDENT_AUDIT.txt}. No claim is made
that the unresolved sentences are open in the entire literature.

\subsection{No-go results for repairing the model}

Increasing $|I|$ alone cannot produce BPI while retaining the absorption
probes, all valid three-copy mergers, root-fixing transport, local
decision, and Cohen conditions mentioning fewer than $\mu$ labels.
Those ingredients already prove the six-cycle obstruction. More
abstractly, the literal infinite-product version of finite Ramsey
homogeneity is false: parity relative to representatives modulo finite
coordinate disagreement gives a binary coloring with no monochromatic
rectangle having infinitely many nonsingleton factors.

Two concrete order-based changes were audited. Replacing the full
coordinate group by automorphisms of a saturated dense linear order
adds a supported cut map $A\twoheadrightarrow\mu^+$, while support
locality still forbids $\mu^+\hookrightarrow A$. Thus PP and
$\mathsf{AC}_{\rm WO}$ fail. Alternatively, decorate every entire
diagram with a ground linear extension of $T$, preserving that order
under restriction. The root-pair amalgamation survives, but the
internal profile map surjects onto all ground linear extensions of
$T$, a set of ground size $2^\mu$. Again PP and ordinal-choice fail.
This differs from finite order forcing over $W$, discussed elsewhere
in this report, whose orbit profiles are new and need not be ground.
See \path{ORDERED_SUBGROUP_CUT_QUOTIENT_OBSTRUCTION.txt} and
\path{ORDER_REPAIR_ROOT_AUDIT.txt}.

For restricted palettes $\mathcal C$ of orders on $T$, closed under
$\vartheta\mapsto\vartheta_t$ with
$x\mathrel{\vartheta_t}y\Longleftrightarrow tx\mathrel\vartheta ty$,
there is no nonempty finite palette, even without requiring extension
of native divisibility. Finite Ramsey produces four probes monotone
up to reversal in every palette order; a five-chart tournament merger
would give a finite set with no least element. For infinite
$|\mathcal C|=\lambda$, the same argument and Erd\H{o}s--Rado give the
necessary bound $\mu\leq2^{2^\lambda}$. It is not an optimal bound.
Palettes of size at most $\mu$ can satisfy pullback closure and every
acyclic finite request. A universal lift can also be constructed for
one prescribed palette, using cross-denominator probe cones sharing
exactly the root. The unresolved issue is closure of one final palette
under universal lifts for all its own witness types, followed by support
descent and quotient classification. No positive-BPI model or universal
impossibility theorem results. See \path{ORDER_PALETTE_ROOT_AUDIT.txt}.

Ransom's 2026 dissertation gives an exact forcing-aware criterion for
BPI (Definition 4.3.1 and Theorem 4.3.3, pp.~93--99), involving dense
sets, bases for the normal filter, orbit colorings, and compatibility
below all common extensions of finitely many translated conditions.
The criterion concerns the full Cohen--diagram product. A Ramsey
property of the abstract monoid alone does not verify it. Its
finite-support examples do not automatically cover these $\mu$-sized
supports. The six-cycle theorem now rules out a successful verification
for the unchanged model. The earlier precise compatibility obligations
remain documented in \path{SYMMETRIC_FILTER_PHASE3.txt} and
\path{SYMMETRIC_FILTER_PHASE4_LARGE_I.txt}.

\subsection{Source corrections, attribution, and remaining boundaries}

The implication $\mathsf{AC}_{\rm WO}\Rightarrow\mathsf{DC}$ is valid
in pure ZF: Jech, Theorem 8.2, pp.~121--122, attributes it to Jensen.
An early apparent contrary table entry was an erroneous HTML
reconstruction, corrected by inspecting the actual primary PDF and
its prose. Here ordinal-choice means choice for well-orderable
\emph{index families}, not arbitrary families of well-orderable
members. The verified chain is
$\mathsf{PP}\Rightarrow\mathsf{AC}_{\rm WO}\Rightarrow\mathsf{DC}$;
it does not yield $\mathsf{DC}_{\omega_1}$.

A different source conflict remains recorded rather than resolved by
recency. Pincus's original announcement and publisher summary, Ransom's
2025 paper, and Asper\'o--Karagila's 2020 manuscript explicitly attribute
models of $\mathsf{ZF}+\mathsf{DC}+\mathsf{BPI}+\neg\mathsf{AC}$ to
Pincus. Kowalczyk--Zapletal, arXiv:2607.16747v1, instead calls the
corresponding implication to AC open; its actual theorem establishes
a restricted compactness fragment. The complete relevant Pincus proof
was not retrieved in this investigation. This discrepancy is not used
as a premise, and no claim that a positive BPI variant would solve
that historical problem is made. See
\path{DC_BPI_SOURCE_CONFLICT.txt}.

The bounded novelty check found no BPI or order-extension statement in
the twelve text files of the pinned PP manuscript and no prior assertion
of this particular PP/FB/NDS countermodel to BPI in the sources checked.
That is limited search evidence, not a priority certificate. The
classical compactness implications, symmetry-cycle method, Ramsey
theorems and coloring constructions retain their established
attributions. The source construction, not this report, proves the
foundational PP consistency theorem. The additional results here are
model consequences and explicitly audited adaptations.

\paragraph{Primary references used in this section.}
Stable source links are:
\href{https://gwern.net/doc/math/1973-jech-theaxiomofchoice.pdf}{Jech (1973)};
\href{https://doi.org/10.2307/2586759}{Felgner--Truss (1999), OEP and BPI};
\href{https://academic.calliope.mx/files/papers/2412.13785.pdf}{Ryan-Smith, Proposition 3.13};
\href{https://arxiv.org/abs/2511.21684v1}{Ransom (2025)} and
\href{https://escholarship.org/uc/item/0b35t7sd}{Ransom's 2026 dissertation};
\href{https://doi.org/10.1215/00294527-2024-0020}{Tachtsis (2024)} and
\href{https://arxiv.org/abs/2310.09643v3}{Blass--Kulshreshtha, version 3};
\href{https://arxiv.org/pdf/math/9311207}{Baumgartner--Hajnal--Todorcevic,
Theorem 2.1};
\href{https://www.renyi.hu/~p_erdos/1956-02.pdf}{Erd\H{o}s--Rado (1956)};
\href{https://users.renyi.hu/~p_erdos/1965-14.pdf}{Erd\H{o}s--Hajnal--Rado
(1965), Lemma 5A}; and
\href{https://www.jams.jp/notice/scmj/56-2.html}{Cowen--Hechler--Mihok (2002)}.
All named local audit files belong to the accompanying
\path{evidence/original_model/} archive directory (mapped from the original \path{BPI_RPP_five_hour_research/} record). Their original
intermediate statuses are preserved; the later independent audits and
this consolidation specify which claims supersede them.

\Needspace{12\baselineskip}\section{The ordinary extension by finite linear orders}
\label{sec:complete-new-model}

This section reports the second model studied in the workspace.  Starting
with the OpenAI symmetric model $W$, we force a linear extension of its
native partial order and obtain $N=W[G]$.  The established conclusions
include
\[
 N\models\mathsf{ZF}+\mathsf{AC}_{\mathrm{WO}}+\mathsf{DC}
       +\neg\mathsf{AC}+\neg\mathsf{DC}_{\omega_1},
 \qquad \mathsf{SVC}(A)+\mathsf{SVC}(\mathbb R_N),
\]
where $\mathbb R_N=\mathcal P(\omega)^N$.  All old cardinal comparisons
survive, but the new real set is not well-orderable.  Whether $N$ satisfies
\emph{global} PP, NDS, FB, or BPI remains unresolved.  In particular, this
section does not present a model of BPI together with failure of Choice.

Here $X\approx Y$ means equipotence, $X\preccurlyeq Y$ means the existence
of an injection, and ``old'' means belonging to $W$, rather than merely
to the earlier ZFC construction ground $V$.  A result described as
audited has a written mathematical proof and a separate workspace
check; this is neither external peer review nor a formal verification
of the whole forcing construction.  The proofs below are summaries of
the saved arguments.  The audit register at the end supplies their
exact filenames.

\subsection{Construction and the chain conditions}

The starting model has a non-well-orderable set $A$ and a ground monoid
$T$ of well-orderable cardinality $\mu$.  There is an internal free right
action $(a,t)\mapsto a\cdot t$ of $T$ on $A$.  ``Free'' here means that
$t\mapsto a\cdot t$ is injective for each fixed $a$; it does not mean
that right multiplication $a\mapsto a\cdot t$ is injective.  Define
\[
 a\leq_T b\quad\Longleftrightarrow\quad
       (\exists t\in T)\ b=a\cdot t.
\]
This is a partial order.  Its principal upset $aT$ is called a
\emph{chart}; it is an old well-orderable set of cardinality $\mu$.
Every old well-orderable subset of $A$ lies in one chart.  The results
used from $W$ include
\[
 A^2\approx A,\qquad A\times\mu\approx A,\qquad
 (A^\omega)^W\approx A,\qquad h^W(A)=\mu^+.
\]
Every old set has a normal form $(A\times\alpha)\sqcup\beta$, with
ordinal coefficients.  Here $h(X)$ is the least ordinal not injecting
into $X$; $w(A)=\mu$ means that $\mu$ is the largest well-orderable
cardinal represented by a subset of $A$.

Let $\mathbb P_{\mathrm{ext}}$ consist of finite strict linear orders on
finite subsets of $A$ which extend the strict part $<_T$ of the native order on their domains.
A stronger condition extends the domain and the finite order.  The
generic union $<_L$ is a strict linear order on all of $A$ extending $<_T$; we also write $L$ for this order.
Set
\[
              N=W[G],\qquad G\subseteq\mathbb P_{\mathrm{ext}}.
\]
This is ordinary set forcing over the choiceless model $W$.  Its
construction and its preservation properties require separate proofs.

\paragraph{Audited chain-condition theorem.}
Every antichain in $\mathbb P_{\mathrm{ext}}$ is countable, and every
predense subset has a countable predense subset.  We call these
properties $\mathrm{CCC}_2$ and $\mathrm{CCC}_3$, respectively.  The
second is not inferred from the first in arbitrary ZF models.

The local proof uses the source monoid's height $\rho:T\to(0,1)$,
positive costs $\epsilon_t$, and conditional joins.  Two conditions
$(x_0<\cdots<x_{n-1})$ and $(y_0<\cdots<y_{m-1})$ are incompatible
exactly when a reversed rectangle occurs:
\[
 i<j,\ l<k,\qquad x_j\leq_T y_l,
                         \quad y_k\leq_T x_i.
\]
After bounding the finite lengths and bounding the costs away from zero,
finite Ramsey theory turns a sufficiently large antichain into a
triangular pattern.  Successive finite joins would then increase
$\rho$ by more than a fixed positive amount at every step, contradicting
$0<\rho<1$.  This proves a local $\sigma$-bounded chain condition on a
chart.  An alleged $\omega_1$-indexed global antichain localizes to one
chart.  In $W$ every uncountable set contains $\omega_1$, so this excludes
all uncountable antichains.

Each chart subforcing is regular in $\mathbb P_{\mathrm{ext}}$.
A finite condition has a finite reduction to the chart: its constraints
are projected to the least native successors in that chart, using
conditional joins.  The stronger $\mathrm{CCC}_3$ proof then uses the
actual source construction's petal-separation lemma.  If a predense set
has supporting diagram $p$, choose a larger diagram $q$ with $q_e=p$.
Its part inside the chart of $q$ is globally predense.  A greedy
well-ordered construction in that chart produces a countable predense
subfamily.  The source-specific petal and gate calculations are material
hypotheses of this argument [N1].

\subsection{Preservation and the remaining failures of Choice}

\paragraph{Audited preservation theorem.}
The extension preserves all ordinal cardinals and cofinalities, the
Hartogs number of every old set, and the non-well-orderability of every
old non-well-orderable set.  More precisely, every new well-orderable
subset of an old set $X$ is contained in an \emph{old} well-orderable
subset of $X$ [N2].

For the proof let $\mathbb B=\operatorname{RO}^W(\mathbb P_{\mathrm{ext}})$.
Its antichains are countable: an uncountable Boolean antichain would
contain an $\omega_1$-indexed subfamily, and old
$\mathsf{AC}_{\mathrm{WO}}$ would select incompatible forcing conditions
below its members.  If a name is forced below $b_0$ to be a function
$f:X\to Y$ between old sets, put
\[
 g(x)=\{y\in Y:b_0\wedge\Vert\dot f(\check x)=\check y\Vert>0\}.
\]
For fixed $x$ the nonzero Boolean values are disjoint.  Thus $g(x)$ is
countable in $W$, and $g$ is one canonically defined old function with
$f(x)\in g(x)$.  For an ordinal domain, old
$\mathsf{AC}_{\mathrm{WO}}$ well-orders the union of these possible-value
sets.  This gives the asserted covering property, from which the
cardinal, cofinality and Hartogs conclusions follow.  In particular
\[
                w^N(A)=\mu,\qquad h^N(A)=\mu^+,
                \qquad N\models\neg\mathsf{AC}.
\]
The complete-Boolean-algebra covering argument is related to the
choiceless chain-condition analysis of Karagila--Schweber; the saved
proof checks the extra hypotheses needed here, instead of importing
ZFC preservation wholesale.\footnote{A.~Karagila and N.~Schweber,
``Choiceless chain conditions,'' \emph{European Journal of Mathematics}
8, Supplement 2 (2022), S393--S410,
\url{https://doi.org/10.1007/s40879-022-00564-2}.}

The stronger chain condition separately preserves
$\mathsf{AC}_{\mathrm{WO}}$.  For each member of a new ordinal-indexed
family, use a countable predense set of conditions carrying witness
names.  Collection first bounds those names in one old set.  Old
ordinal-indexed choice selects the countable lists and witnesses;
in the generic extension, the first successful condition supplies the
chosen element.  This uses local witness density, not the choiceless
maximum principle.  Consequently $N$ also satisfies ordinary DC,
by the established implication $\mathsf{AC}_{\mathrm{WO}}\Rightarrow
\mathsf{DC}$ [N2].

Every new ordinal sequence of $A$-points is contained in one old chart.
It therefore has a strict common native predecessor.  A descending
$\omega_1$-sequence would produce, in that chart, a strictly decreasing
$\omega_1$-sequence of real heights, which is impossible.  The resulting
countably closed tree of descending native sequences witnesses
$\neg\mathsf{DC}_{\omega_1}$.  Also $L$ has no well-orderable coinitial
subset.  Its reverse witnesses failure of the Cofinality Principle,
the assertion that every linear order has a well-orderable cofinal
subset.  These statements concern new subsets as well as old ones.

The old Boolean algebra $\mathbb B$ has cardinal $A$, is complete,
atomless and ccc in $W$, and has no ultrafilter there: an ultrafilter
would evaluate its comparison values to a forbidden native linear
extension.  The generic extension supplies an ultrafilter on this
\emph{old} algebra.  Its old completeness need not extend to all new
families of Boolean elements.  Repairing this particular obstruction
does not establish BPI for all Boolean algebras of $N$ [N3].

\subsection{The new real set and the unresolved evaluation quotient}

\paragraph{Audited real-cardinal theorem.}
In $N$ there are maps
\[
            A\lhook\joinrel\longrightarrow\mathbb R_N,
       \qquad A\twoheadrightarrow\mathbb R_N.
\]
The real set is not well-orderable and satisfies
\[
 w(\mathbb R_N)=\mu,\qquad h(\mathbb R_N)=\mu^+,
 \qquad \ell(\mathbb R_N)=\mu^+,
 \qquad (A^\omega)^N\approx\mathbb R_N.
\]
For the original parameters $\mu=\aleph_2$, this gives the non-well-orderable
real set Hartogs number $\aleph_3$; parameter variants use their preserved
$\mu$.  Here $\ell$ is the least positive ordinal not obtainable as a
surjective image; the harmless empty-set convention does not affect
the displayed infinite value [N4].

For the injection, fix selected probes $m_n$ and encode the order
pattern of $\{a\cdot m_n:n<\omega\}$ as a real $r_a$.
The probe geometry implies that the sets at two distinct centers
intersect in at most one point.  A finite-order swap can change a
comparison at one center while keeping the corresponding comparison
at the other unchanged.  These separating extensions are dense, so
$a\mapsto r_a$ is injective.

For the surjection, $\mathrm{CCC}_3$ expresses every old Boolean value
as a countable join of forcing conditions.  Since
$|\mathbb P_{\mathrm{ext}}|=|A|$ and $(A^\omega)^W\approx A$, one old
$A$-sized set codes all old sequences of Boolean values.  Evaluation
under the generic ultrafilter covers every new subset of $\omega$.
This produces a surjection, not an injection in the opposite direction.

If $\mu^+$ injected into $\mathbb R_N$, new
$\mathsf{AC}_{\mathrm{WO}}$ would choose preimages under
$A\twoheadrightarrow\mathbb R_N$, contradicting the preserved Hartogs
number of $A$.  This proves the upper threshold; the profile injection
gives the lower bound.  Finally,
$\mathbb R_N\preccurlyeq A^\omega\preccurlyeq
\mathbb R_N^\omega\approx\mathbb R_N$ proves the countable-power
identity.  The left-hand power here is computed in $N$, not in $W$.

The central unresolved comparison is
\[
                         \mathbb R_N\preccurlyeq A\ ?
\]
An injection and a surjection from $A$ to the same target do not
settle this in ZF.  Least choices of ground code \emph{types} do not
choose centers in the remaining non-well-orderable fibers.  Even
$\mathbb R_N\approx A$, if eventually proved, would settle one
quotient instance and would not alone prove global PP.

\subsection{Local real fields and exact localization}

Put $Q=\mathbb P_{\mathrm{ext}}(T)^V$.  For each $a\in A$, the chart
map $t\mapsto a\cdot t$ pulls the generic order back to a $Q$-generic
filter $K_a$.  Chart regularity gives genericity over $W$, hence over
$V$; it does not justify genericity over the external Cohen universe
$V[F]$.  Fix in $V$ a $\mu$-enumerated set $\Gamma$ of all countable
$Q$-names coding reals, and define internally
\[
 R_a=\{\tau^{K_a}:\tau\in\Gamma\}.
\]
Externally this is $\mathbb R_N\cap V[K_a]$.  The internal definition
uses a set of names, not a predicate for the ground universe.

\paragraph{Audited local-field theorem.}
For every $a$, $|R_a|=\mu$.  Every new well-orderable subset of
$\mathbb R_N$ is contained in some $R_a$.  Every new subset of an
ordinal belongs to $W[K_a]$ for some $a$.  Moreover
\[
 R_a\subseteq R_b\quad\Longleftrightarrow\quad a\in bT,
 \qquad
 r_a\in R_b\quad\Longleftrightarrow\quad a\in bT.
\]
Indeed the stronger exclusion $r_a\notin W[K_b]$ holds whenever
$a\notin bT$ [N5].

The size bound comes from the ground enumeration of $\Gamma$; the
$\mu$ profiles $r_{a\cdot t}$ give equality.  For localization, take
the Boolean bits of an ordinal-indexed real family and choose their
countable forcing codes using old ordinal-indexed choice.  All finite
point domains form one old well-orderable subset of $A$, hence lie
in one chart.  Each individual real then has a countable sequence of
ground $Q$-conditions.  The original model adds no countable sequences
of ground elements, so that code belongs to $V$ and evaluates in
$V[K_a]$.  The full uncountable array need only belong to $W$; this
explains the weaker $W[K_a]$ conclusion for arbitrary ordinal subsets.

The inclusion-reflection proof uses a quotient-forcing argument.  If
$a\notin bT$, infinitely many of the probes at $a$ have identical
native comparison profiles to the entire chart $bT$.  Below any finite
quotient condition, two remain interchangeable while all chart-order
information is fixed.  Both orientations can still be forced, so the
corresponding bit of $r_a$ is new over $W[K_b]$.

Two distinctions are essential.  A well-orderable real family is
\emph{contained in} $R_a$; the family itself need not belong to
$V[K_a]$.  Also proper field inclusions do not give strict cardinal
inequalities: all the $R_a$ have the same cardinal $\mu$.  Every
well-orderable family of these fields is contained in one strictly
larger field, by taking a strict native predecessor of its centers.
Thus there is no well-orderable cofinal subfamily under inclusion.

\subsection{All comparisons between old sets are preserved}

For any old $X,Y$, the existence in $N$ of an injection, bijection, or
surjection between them is equivalent to its existence in $W$.
Every new surjection between old sets has the PP conclusion witnessed
by an \emph{old} injection in the reverse direction.  Moreover every
new set-family consisting of old sets has a finite coinitial cardinal
basis [N6].

Here is the comparison invariant.  An old well-orderable set of
cardinality $\nu$ has record $(0,\nu)$.  Normalize a non-well-orderable
old set as
\[
 X\approx(A\times\alpha)\sqcup\beta,\qquad
 \alpha\geq\mu,\quad \beta=0\text{ or }\beta>\alpha,
\]
with ordinal cardinals as coefficients, and give it record
$(\alpha,\varepsilon)$, where $\varepsilon=\max(\alpha,\beta)$.
These records are unique and injection comparison is componentwise,
both in $W$ and in $N$.

For example, an injection from record $(\alpha,\varepsilon)$ into
$(\gamma,\zeta)$ gives $\varepsilon\leq\zeta$ by Hartogs preservation.
Each of the $\alpha$ source copies of $A$ must meet the target's
$A$-part; otherwise that copy would become well-orderable.  New
ordinal-indexed choice selects one point from each such disjoint
intersection, giving $\alpha\leq\gamma$.  Conversely the two
inequalities already construct an old injection.  A similar
preimage selection proves preservation of surjections.

Every set of ordinal pairs has a finite coinitial subset under the
componentwise order: repeatedly choose a least first coordinate and
then a least second coordinate not already dominated.  Continuing
infinitely would strictly decrease the second coordinate infinitely.
Applying this to the unique old records proves the relative FB claim.
For a new family of old sets, first bound their ranks inside one old
rank segment, so its old record assignment is a set function.

Thus an infinite new cardinal antichain or descending chain cannot
consist solely of old cardinalities.  This is not global FB or NDS
for all new sets.  If $\mathbb R_N\not\approx A$, its cardinality is
new: $A$ is the only old non-well-orderable cardinality with
$w=\mu$.

\subsection{Explicit nonsplitting and additional choice bounds}

For every $t\neq1$, the old surjection $f_t:A\to A$,
$f_t(x)=x\cdot t$, remains without a right inverse.  More strongly,
there is no assignment of a nonempty countable subset of each fiber,
even in $N$ [N7].  The old proof says that any countable-valued
assignment on $A$ has some value contained in the forward chart of
its argument.  A nonempty value consisting of predecessors would
then contradict conicality of $T$.  For a new assignment,
$\mathrm{CCC}_3$ supplies a canonical old countable cover of every
possible value, reducing the contradiction to that old statement.

Transport $f_t$ along $a\mapsto r_a$ and extend it by the identity
outside that image.  This gives a surjection
$g_t:\mathbb R_N\to\mathbb R_N$ without a right inverse, or even a
countable multisection.  Consequently choice fails for real-indexed
families of nonempty sets of reals.  This is not a failure of PP:
PP does not split each specified surjection, and these two sets are
already equipotent.

Ordinary forcing preserves $\mathsf{SVC}(A)$.  The profile injection
gives a retraction $\mathbb R_N\twoheadrightarrow A$, hence
$\mathsf{SVC}(\mathbb R_N)$.  Coding
$\mathbb R_N\times\eta$ into a power set of an ordinal yields
\[
             \mathsf{KWP}^{*}_1,\qquad \mathsf{KWP}_2
             \quad\hbox{in }N.
\]
Here $\mathsf{KWP}_n$ means that every set injects into an $n$-fold
power set of an ordinal; the starred version uses partial surjections.
Inverse-image fibers give $\mathsf{KWP}^{*}_1\Rightarrow
\mathsf{KWP}_2$.  These are the established bounds, not
$\mathsf{KWP}_1$, the Ordering Principle, or global finite-choice
preservation.\footnote{A.~Karagila and J.~Schilhan,
``Intermediate models and Kinna--Wagner Principles,''
\emph{Proceedings of the American Mathematical Society} 154 (2026),
393--403, \url{https://doi.org/10.1090/proc/17425}; see Section 3.2
and Proposition 4.1 of \url{https://arxiv.org/html/2409.07352v2}.}

There are also definite negative continuum-function conclusions.
Write ICF for
$\mathcal P(X)\approx\mathcal P(Y)\Rightarrow X\approx Y$,
and HCF for reflection of injection comparison by $\mathcal P$.
ZF proves $\mathsf{ICF}+\mathsf{KWP}_{n+1}\Rightarrow
\mathsf{KWP}_n$: apply the latter to $\mathcal P(X)$, obtaining
$\mathcal P(X)\preccurlyeq\mathcal P(B)$ with
$B=\mathcal P^n(\theta)$ and $\theta$ infinite.  Since $B$ absorbs
two copies,
$\mathcal P(X\sqcup B)\approx\mathcal P(B)$, so ICF gives
$X\preccurlyeq B$.  Finite descent reaches AC.  Thus the proved
$\mathsf{KWP}_2+\neg\mathsf{AC}$ package implies
$\neg\mathsf{ICF}+\neg\mathsf{HCF}$ in $N$ [N8].

\subsection{Initial segments, final segments, and open intervals}

For $a<_L b$ put
\[
 I_a=\{x:x<_L a\},\qquad J_a=\{x:a<_L x\},\qquad
 (a,b)_L=\{x:a<_L x<_L b\}.
\]
The final audited result is
\[
                  I_a\approx J_a\approx(a,b)_L\approx A.
\]
The earlier possibility that a final segment might be a proper
square root of $A$ was resolved by an explicit injection, not by
cancelling the established identity $J_a^2\approx A$ [N9].

An initial segment contains the old $A$-sized fiber
$\{x:x\cdot t=a\}$ for any $t\neq1$.  For an interval, choose selected
probes $m_n$ with $m_ne=e$, pairwise cone intersection $eT$, and only
identity common left divisors.  Fix $d$ with $d\cdot e=b$, and set
$F=\{x:x\cdot m_0=d\}\approx A$.  For each $x\in F$ the ray
\[
                    S_x=\{x\cdot m_n:n\geq1\}
\]
lies strictly below $b$.  These rays are disjoint for different $x$:
a shared ray point together with the shared anchor $d$ forces the
common-chart join quotients to divide both $m_0$ and a later $m_n$,
so both quotients are $1$.

In a finite condition, a point below $b$ can dominate at most one
point of $S_x$; dominating two would place it natively above $b$.
A remaining ray point can therefore be inserted immediately before
$b$.  Genericity makes $S_x$ cofinal below $b$.  Its least indexed
point above $a$ injects $F$ into $(a,b)_L$, proving the theorem.
Taking $b=a\cdot e$ also proves the final-segment case.

The whole $L$ has countable cofinality.  A fixed countable sequence
of fresh merger letters with disjoint cones gives, above any center,
an old countable cofinal ray.  Every old non-well-orderable subset of
$A$ is nevertheless $L$-coinitial: a point outside a chart covering
a finite condition can be placed first.  Thus no final segment or
bounded interval contains an old non-well-orderable subset.  The
$A$-sized copies just constructed are new.

All these segments have finite positive powers of cardinal $A$ and
new countable power of cardinal $\mathbb R_N$.  They therefore yield
no strict cardinal descent among themselves.  A further recorded
refinement says that each final segment has an old countable
order-dense subset, using the common-divisor and disjoint-cone
properties of the merger rays.  That refinement has a written proof
in [N9], but its separate independent audit is not part of the saved
segment audit; it is not used in any theorem above.

\subsection{Boundaries of the current classification}

The unresolved global questions in $N$ are PP, the strong quotient
property $Q(A)$, NDS, FB, the assertion that every cardinal antichain
is finite, and BPI.
Also unresolved are $\mathbb R_N\preccurlyeq A$,
$\mathsf{KWP}_1$, the Ordering Principle, and Choice for arbitrary
families of finite sets.  The original model's countable-member choice
consequence has not been transferred to $N$.  Adding a native linear
extension does not establish the Order Extension Principle for every
partial order, and neither it nor the local fields settles BPI.

Two bounded deductions sharpen the outstanding work.  First, a
hypothetical map $j:\mathbb R_N\to A$ has a support root $d$ such
that $j[R_a]\subseteq aT$ whenever $a\leq_T d$.  This produces a
coherence requirement among local $\mu$-sized injections, rather than
a cardinal contradiction.  Second, every directed ideal of $T$ in
$N$ has a canonically chosen countable cofinal sequence: select least
points in the old $T$-wellorder with heights approaching the supremum,
then take finite joins.  The positive-cost inequality rules out an
omitted divisor.  Hence all these ideals inject into
$\mu^\omega\approx\mathbb R_N$ [N10].  These are saved local
deductions; neither supplies the missing global coherent selector.

The original quotient proof explicitly requires support descent and
coherent matching across charts.  New availability ideals need not
be ground objects, and arbitrary ordered three-chart amalgams can
cycle.  A bound on local names or ideals by $\mathbb R_N$ is therefore
not a proof of an injective presentation theorem.  Likewise the
Truss-model comparison in Ryan-Smith's work shows why
$\mathsf{AC}_{\mathrm{WO}}+\mathsf{SVC}(\mathbb R)$ alone must not be
treated as PP.\footnote{C.~Ryan-Smith, ``Local reflections of choice,''
\emph{Acta Mathematica Hungarica} 176 (2025), 244--257,
\url{https://doi.org/10.1007/s10474-025-01533-3}.
The relevant comparison is Proposition 3.19 in the published version,
numbered 3.18 in the author-hosted preprint.  It concerns a different
model; it is not a negative PP theorem about $N$.}

\subsection*{Saved proof and status register for this section}

All filenames below are in the public archive directory
\path{evidence/original_model/}, mapped from the workspace record
\path{outputs/BPI_RPP_five_hour_research/}.  The source references are
the pinned OpenAI construction in \path{work/openai_math/}: in
particular \path{forcing.tex} (diagram amalgamation, local decision,
ordinal choice and no countable ground sequences), \path{monoid.tex}
(probe gates, merger normal forms and the construction schedule), and
\path{quotients.tex} (availability ideals and coherent quotient
matching).  The external construction is credited to
\url{https://github.com/openai/math}; the extensions and deductions
reported here must be distinguished from statements made in that
manuscript.

\begingroup\small
\noindent\textbf{[N1] Chain conditions and chart reduction, audited.}
\path{FINITE_LINEAR_EXTENSION_FORCING_CCC.txt};
\path{FINITE_LINEAR_EXTENSION_CCC_INDEPENDENT_AUDIT.txt};
\path{FINITE_ORDER_FORCING_CCC3_CHART_REDUCTION.txt};
\path{FINITE_ORDER_FORCING_CCC3_INDEPENDENT_AUDIT.txt}.

\noindent\textbf{[N2] Preservation, audited.}
\path{BOOLEAN_COMPLETION_COVERING_AND_HARTOGS_PRESERVATION_AUDIT.txt};
\path{CCC3_PRESERVES_AC_WO_CONDITIONAL_AUDIT.txt};
\path{AC_WO_DC_PRIMARY_SOURCE_RECONCILIATION.txt}.
The ``conditional'' filename records an earlier checkpoint: its
$\mathrm{CCC}_3$ hypothesis was subsequently verified in [N1].

\noindent\textbf{[N3] The old complete Boolean algebra, audited.}
\path{COMPLETE_CCC_BOOLEAN_ALGEBRA_EXACT_CARDINAL_A_AUDIT.txt};
\path{BOOLEAN_ALGEBRA_ROOT_AUDIT.txt}.

\noindent\textbf{[N4] Real maps and thresholds, audited.}
\path{GENERIC_ORDER_PROFILES_INJECT_A_INTO_REALS.txt};
\path{GENERIC_REAL_PROFILE_INDEPENDENT_AUDIT_AND_QUOTIENT_TEST.txt};
\path{BOOLEAN_CODES_FOR_NEW_REALS_ROOT_DRAFT.txt};
\path{BOOLEAN_CODES_REAL_QUOTIENT_INDEPENDENT_AUDIT.txt}.
The retained ``draft'' filename is accompanied by the separate audit.

\noindent\textbf{[N5] Local fields and localization, audited.}
\path{LOCAL_REAL_FIELDS_REFLECT_NATIVE_ORDER.txt};
\path{LOCAL_REAL_FIELDS_INDEPENDENT_AUDIT.txt};
\path{LOCAL_ZFC_REAL_COVERS_AND_ORDINAL_LOCALIZATION_AUDIT.txt}.

\noindent\textbf{[N6] Old comparisons and relative finite bases.}
\path{OLD_CARDINAL_COMPARISONS_AND_RELATIVE_FINITE_BASIS.txt};
\path{OLD_SURJECTIVE_COMPARISONS_AND_RELATIVE_PP.txt}.
These are mathematical deductions from the audited preservation and
normal-form results; the surjective comparison has a recorded second
reconstruction.

\noindent\textbf{[N7] Nonsplitting and multisections, audited.}
\path{COUNTABLE_MULTISECTIONS_AND_NONSPLITTING_PRESERVATION.txt};
\path{COUNTABLE_MULTISELECTOR_INDEPENDENT_AUDIT.txt}.

\noindent\textbf{[N8] Choice bounds and continuum-function failure.}
\path{NEW_MODEL_KWP_OP_FINITE_CHOICE_BOUNDARY.txt};
\path{FINITE_KWP_ICF_COLLAPSE_INDEPENDENT_AUDIT.txt}.

\noindent\textbf{[N9] Resolved segment cardinalities, audited.}
\path{ANCHORED_RAYS_RESOLVE_TAIL_AND_INTERVAL_CARDINALS.txt};
\path{GENERIC_ORDER_SEGMENTS_INDEPENDENT_AUDIT.txt}.
The preliminary file
\path{GENERIC_LINEAR_ORDER_SEGMENTS_AND_SQUARE_ROOTS.txt}
now carries a supersession notice.  Its open tail alternative is not
a current unresolved question.  The separability refinement in the
later note remains separately flagged above.

\noindent\textbf{[N10] Bounded further deductions and remaining gaps.}
\path{REAL_TO_LABEL_INJECTION_CHART_LOCALITY_AUDIT.txt};
\path{NEW_AVAILABILITY_IDEALS_REAL_CODE_BOUND.txt};
\path{TRUSS_QUOTIENT_COMPARISON_AFTER_ORDER_FORCING.txt}.
These do not assert a completed new quotient classification.
\endgroup

\clearpage\section{Earlier permutation-model results and construction obstructions}
\label{sec:earlier}
These results belong to the earlier ordinary Fraenkel--Mostowski (FM) search. They are retained because a failed PP construction may still satisfy useful restricted principles, and because the obstructions explain why merely modifying a group or filter is difficult. The ground universe has atoms and choice. The action is faithful, the normal filter is effective, and ``open'' means allowed by that filter. All densities, indices and multiplicities in this section are external ground cardinals. Claims about specific constructions are conditional on their recorded normal-form and extension lemmas; they are not claims about every permutation model.

\subsection{Exact and sufficient group criteria}
\paragraph{The exact PP criterion.}
For every open $K$, there must be one open $J\leq K$ such that for every pair of open subgroups $L\leq H\leq K$ there is a $J$-equivariant injection
\[
\operatorname{Res}_J(K/H)\inj\operatorname{Res}_J(K/L).
\tag{P}
\]
The same $J$ must work for all pairs. Matching requires the exact conjugacy type of each point stabilizer and enough target orbits of that type. A fixed point, an inclusion between stabilizers, or individual repairs with unbounded supports do not suffice. The criterion is hereditary to open subgroups; passing to a dense subgroup with the induced topology preserves it.

\paragraph{Common compression.}
The group formulation of $\ACWO$ used here is
\[
\forall K\text{ open}\ \exists Q\leq K\text{ open}\
\forall L\leq K\text{ open}\ \exists k\in K\quad k^{-1}Qk\leq L.
\tag{C}
\]
In a faithful effective action, failure of choice is equivalent to the identity subgroup not being open. A coset object $K/H$ is internally well-orderable precisely when its normal core in $K$ is open: a supported well-order has trivial permutation action on its domain, and an open core conversely supports a ground well-order.

\paragraph{Uniform non-well-orderable profiles.}
A sufficient additional condition is that, after one open refinement $J\leq K$, all $K/H$ with non-open core become isomorphic as $J$-sets. Combined with (C), well-orderable targets are handled by simultaneous sections and the remaining targets by these common bijections. This implies (P). It also gives a two-prototype decomposition and hence NDS. No nondiscrete realization of this uniform condition was obtained.

\paragraph{Finite prototypes.}
More generally, suppose that for each open $K$ there are an open $J\leq K$ and finitely many continuous $J$-sets $U_1,\ldots,U_r$ such that every restricted $K/H$ is a disjoint union of ground-cardinal multiples of these same prototypes. Decomposing a supported sequence into such orbits gives coefficient vectors in a finite product of well-orders. An increasing pair of vectors gives a supported reverse comparison, excluding strict descent. This proves NDS and finite antichains, without a uniform finite bound on their sizes. The condition also implies (C), by fixing a point of each nonempty prototype. Its realization remains unresolved.

\paragraph{Relative point completeness (RPC).}
For a transitive atom action with cofinal point fixers, (P) requires a refinement below which every open subgroup is the exact stabilizer of an atom. Equivalently, at a suitable root $b$, every open $H\leq G_b$ has the form $\operatorname{Fix}(b,a)$. In particular such a subgroup is determined by its fixed atoms; arbitrary increasing unions must retain this property. In this setting, (C), RPC and the dichotomy that every internal atom subset is well-orderable or equipotent to the atom set imply full PP. This sufficient criterion has no realized nondiscrete example in the record. A multi-HNN construction realizes (C) and RPC with failure of AC, but fails the dichotomy, PP and NDS.

\subsection{Necessary topology and a sharp atom-model boundary}
The following restrictions have saved proofs and separate checks.
\begin{enumerate}
\item Under (C), the intersection $R(K)$ of the open normal subgroups of $K$ is open and satisfies $R(R(K))=R(K)$. It has no proper open subgroup of finite index. A commensurated open subgroup has open core; a cofinal commensurated neighbourhood base therefore forces discreteness.
\item A nondiscrete (C)-group has no open compact or open Roelcke-precompact subgroup. Cofinally one obtains an open $B\rtimes\langle g\rangle$ with quotient $\mathbb Z$. The contraction argument gives a proper conjugate copy inside a common-compression witness; co-Hopfian, J\'onsson, or malnormal witness designs are consequently unsuitable in the stated setting.
\item A nondiscrete effective FM group satisfying PP has
\[
 d(K)\geq 2^{\aleph_0}\quad\text{for every open }K.
\]
The proof interpolates a rationally dense chain in one open-subgroup interval. Real cuts yield continuum many open intermediate subgroups, each requiring a distinct exact stabilizer in one coset target. A dense subset meets all open cosets, giving the cardinal bound. Thus no separable open subgroup is possible. This excludes Polish effective topologies, not every topology on a countable atom set.
\item If $J$ is a common-compression witness in $K$, then $[K:J]=d(K)$ and the character and weight are bounded by this index (the weight equals $d(K)$ in the recorded setting). Below a PP witness, the number of open subgroups is at most $d(K)$. A finite-prototype witness satisfies the same local capacity bound. No general amplification to $2^{d(K)}$ was proved.
\end{enumerate}

\paragraph{Finite choice in ordinary FM.}
For each fixed $n\geq2$, $\ACWO+\mathsf{AC}_n$ implies $\AC$ in the ordinary FM setting above. Here $\mathsf{AC}_n$ means choice for arbitrary families of $n$-element sets. The proof uses tail-equivalence classes of injective atom sequences. Ordinal-indexed choice supplies a common support for representatives of well-ordered configurations. If that support still moves an atom, a contracting translate creates an exact $n$-cycle of tail classes, contradicting a supported selector. Therefore the support fixes all atoms and is trivial.

Consequently $\ACWO+\mathsf{OP}$ collapses to $\AC$ in ordinary FM. The OpenAI generator $A\subseteq\mathcal P(\kappa)$ has a lexicographic order; its full linearly orderable strong-generator package cannot occur in a nonchoice ordinary FM model. This is an atom-model obstruction, not a contradiction in pure $\ZF$.

If an unordered strong generator did exist in ordinary FM, it could be replaced by a subset of atoms of least external nonchoice size $\lambda$. After an open restriction, its point stabilizers would form a cofinal neighbourhood base, and every further open subgroup would have density exactly $\lambda$. Thus $\lambda\geq2^{\aleph_0}$. These constraints do not construct such a generator or prove that it is impossible.

\subsection{Germ and monoid completions: what succeeded locally}
In the imported free germ construction, complete stationary-tail quotients and countable-prefix quotients admit reverse injections into the atom set with one prescribed root support. The proof preserves exact relative stabilizers through triangular extensions and provides enough tagged target orbits. This is a whole-quotient result, not just a local finite test.

Proper-suffix, nonfree whole-act extensions improve this to exact codes for ordered tuples and unordered finite configurations. In the mixed completion, each nontrivial supported atom orbit contains a copy of the whole atom set after one common refinement. Hence every atom subset is well-orderable or equipotent to $A$, and PP and NDS hold when restricted to atom subsets. With the recorded countable-code schedule, $\operatorname{Count}(A)\cong A$ and every surjection from an atom subset with countable fibers has a reverse injection. The argument does not choose enumerations of all fibers simultaneously.

Zero-suffix markers repair the residual stabilizer groups of the proper-suffix construction: a fair schedule kills every final group $U(P_f)$, where $P_f=\{u:uf=f\}$. Basic root fixers then have no proper open normal subgroup. Arbitrary open subgroups may still have nontrivial quotients. Neither this repair nor countable-fiber PP establishes full PP.

\subsection{Explicit counterexamples to the germ candidates}
\paragraph{Free finite-period germs.}
A preserved divisibility invariant says that certain periodic fixed labels meeting a designated cone already lie in that cone. It leaves a target coset with an exact relative stabilizer absent from every possible image point. A stronger direct witness is
\[
\operatorname{Anc}(u)\times M\onto A,\qquad (x,r)\longmapsto xr,
\]
where ancestors of $u$ have no supported reverse injection from $A$. Thus the final imported Section 85 candidate fails PP, including its optional local-group completion. Global absorber resets do not remove the endpoint-specific invariant. This proof does not establish failure of NDS in that candidate.

\paragraph{Proper-suffix-only completion.}
It repairs that ancestor witness but preserves an action-groupoid residual. Tagged loops yield $2^\lambda$ open overgroups of one point fixer, while the relevant target coset has only $\lambda$ points. Exact stabilizer capacity refutes full PP. Its restricted positive results remain valid; NDS is unclassified here.

\paragraph{Mixed zero-suffix completion.}
The residual-group obstruction is repaired, but a long-chain quotient supplies a different counterexample. The scheduled operations produce decreasing principal ideals $h_\alpha M$ with
\[
\bigcap_{\alpha<\omega_1}h_\alpha M=aM.
\]
The relation $xEy$ iff $xh_\alpha=yh_\alpha$ for some $\alpha$ defines an internal surjection $A\onto A/E$. Its class stabilizer is a proper increasing union with the same fixed atoms as a larger root fixer. A supported reverse injection is impossible. The construction generalizes at the stated regular scheduling boundaries. Countably based repairs preserve this obstruction; arbitrary uncountably based repairs are not excluded. Failure of NDS does not follow from this PP counterexample.

\Needspace{12\baselineskip}
\subsection{Fields, local-motion groups, and symmetric iterations}
\begin{longtable}{>{\raggedright\arraybackslash}p{.26\linewidth}>{\raggedright\arraybackslash}p{.66\linewidth}}
\toprule
Construction or method & Precise outcome recorded\\\midrule\endhead
Countable full-field group & Dense chains of root fields give continuum many distinct cut fixers in an interval whose source enumeration orbit is countable. PP fails; the literal uncountable-chain proof is unnecessary. NDS remains unclassified.\\
Uncountable algebraically closed atom field with countable closed-subfield supports & A proper dense subgroup of a countable field automorphism group has the same fixed atoms as its closure. The enumeration quotient violates exact stabilizer matching, so PP fails.\\
Saturated full separably closed fields & With the recorded named finite $p$-basis and sufficient saturation/homogeneity, every fixed finite or infinite support bound is excluded for PP. Kummer root capacity, dense restriction groups, or a binary-profile surjection handle the separate size cases. Arbitrary normal filters and restricted nonsaturated groups are outside this conclusion.\\
General saturated structures & Countable language/base with finite supports fails $\ACWO$ under the stated saturation hypotheses. For language/base of size $\leq\lambda$ and all $\leq\lambda$ supports, order-indiscernible restriction groups violate PP capacity.\\
Increasing homeomorphisms of a ray & With initial-segment fixers, (C) and nonchoice hold and every open subgroup has continuum density. The endpoint quotient refutes PP; separate prime-localization coset tails refute NDS. First countability alone is not enough.\\
Finite piecewise-affine interval exchanges & Bounded-interval supports give (C) and nonchoice. Parameter subspaces give too many overgroups for PP; prime-cycle tails independently refute NDS. Removing an atom order does not repair the problem.\\
Continuous local-motion repairs & On the stated second-countable patches, a countable moving subgroup has the same fixed atoms as a larger patch group, giving the quotient obstruction. In the affine example, compact motion is dense, so a Hausdorff quotient killing it is trivial.\\
Full-symmetric enlargement & A genuine non-open, trace-preserving extension supplies exact old stabilizer types and repairs enlarged old quotient diagrams. The larger group creates new PP and NDS obstructions.\\
Iterated full-symmetric absorption & At cofinality above the support bound the limit is dense in a full symmetric group and fails PP and NDS. At smaller cofinality, a reservoir diagonal defeats (C) and gives a well-orderable-target PP obstruction. Small-cofinality NDS is not settled.\\
\bottomrule
\end{longtable}

A separate algebraic lemma is useful in the field exclusions: if an uncountable field $C$ lies in a separably closed field with a transcendental $t$ over $C$, the latter requires at least $|C|$ field generators over $C$. Roots of $t-a$ for distinct $a\in C$, after the stated finite cyclotomic adjustment, give the Kummer degree bound. The conclusion concerns literal field generation, not a closure operation that already includes separable closure.

\subsection{Coherent coding is possible, but the constructed models fail}
The forest constructions realize actual equivariant coding maps, so their positive part should not be confused with a merely proposed coding scheme.
\begin{enumerate}
\item A binary forest realizes $[A]^2\cong A$ equivariantly; fixed finite arity $m$ similarly realizes $[A]^m\cong A$. Appropriate infinite support filters yield (C) and nonchoice. A fresh cyclic configuration of prime order greater than $m$ refutes PP, and tails over distinct such primes refute NDS.
\item A general finite-carrier theorem removes any bound on signature size: if every atom has an actual finite pointwise support under an embedded full basis permutation group, the quotient of injective sequences by finite coordinate permutations is a PP counterexample under the stated infinite support bound. Its fibers are countable. A merely finite carrier label is not enough for this theorem.
\item A countably branching forest realizes codes for all countable subsets and exact sequence codes. It repairs the previous finite-reordering quotient and finite prime-cycle tails. Nevertheless the fixed-support construction has too many open overgroups for PP.
\item NDS in that countable-code model fails by a different mechanism: hidden index-$p$ subgroups of compact auxiliary product groups have the same fixed atoms as the whole product. The exact coded finite covers give strictly descending tails over distinct primes. Higher fixed infinite arities and the singular-bound version admit the corresponding capacity argument.
\item At singular $\lambda$, the recorded injective $<\lambda$ codes bootstrap to exact $\lambda$-sequence codes, causing support collapse and the previous obstructions. At regular uncountable limit $\lambda$, the particular forest fails common compression. A general theorem excluding all unbounded-arity or huge-signature constructions was not proved.
\end{enumerate}

\paragraph{Proof records.} The supporting archive includes the original follow-up mathematical notes under \path{evidence/earlier_FM/}, including \path{NO_SEPARABLE_OPEN_GROUP_PP_THEOREM.txt}, \path{RELATIVE_POINT_COMPLETENESS_CRITERION.txt}, \path{MIXED_COMPLETION_LONG_CHAIN_PP_COUNTEREXAMPLE.txt}, \path{SATURATED_SCF_FIXED_SUPPORT_BOUND_NO_GO.txt}, \path{COHERENT_COUNTABLE_CODES_AND_CAPACITY_OBSTRUCTION.txt} and \path{HIDDEN_PRIME_COVER_NDS_OBSTRUCTION.txt}. Earlier ``current candidate'' labels in those chronological notes are superseded by the conclusions above.

\subsection{Other imported forcing results and their limits}
The untouched imported report also records restricted pure-$\ZF$ forcing analyses. These were historical source material, not a second independently certified construction in the present consolidation. Their retained conclusions are:
\begin{itemize}
\item Ordinal-target PP is equivalent to $\ACWO$; the first failure of well-orderability in pure $\ZF+\ACWO$ occurs at a power set of an ordinal. Under $\SVC^+(S)$, full PP localizes to $\ACWO$ plus PP between subsets of $S$. Whole-source PP$(S)$ is a separate condition when only $\SVC(S)$ is known.
\item Repairing old real quotients, keeping old reals, or proving PP between old objects does not control new sources. Externally ccc repairs and internally well-orderable repairs have distinct exclusions over the specified Feferman seed.
\item The recorded stabilized repair handles quotients of the full final real set, with dependent choice and $\SVC^+(\mathbb R)$ under its stated assumptions. Its unrestricted ordinal-choice/distributivity and arbitrary-source obligations were not closed. A sparse variant fails PP and the relevant ordinal-choice principle; that does not refute the different unbiased version.
\item Specified Sacks, Silver, Mathias, Laver, Miller and related routes are excluded only under the report's seed and preservation hypotheses. Ordinary splitting and fat splitting are not completely classified there. A no-proper-intermediate theorem retaining the seed does not cover submodels that discard old reals.
\item An amoeba proposal lacks its countable-pattern properness proof; the full-condition product proposal has an explicit choice cone whose density is not proved. These are unfinished constructions, not consistency theorems.
\end{itemize}
Bounded-rank Jech--Sochor transfer alone cannot turn an ordinary-FM statement with unrestricted all-set quantifiers into the desired pure-$\ZF$ theorem. The positive pure model used in the current results is the separately cited OpenAI construction.

\Needspace{12\baselineskip}\section{Verification, novelty and the remaining research programme}
\label{sec:verification}
\subsection{What was checked formally}
The saved isolated build compiled all 75 unchanged pinned \path{PartitionConsistency} proof files with Lean 4.34.1 and the pinned Mathlib. The recorded \texttt{\#print axioms} checks cover the main consistency declarations and the stronger inputs: non-well-orderable quotient/image bijections, ordinal presentation, and ordinal-indexed choice. Only \texttt{propext}, \texttt{Classical.choice}, and \texttt{Quot.sound} occur in those checked declarations. The construction sources were identity-checked. This is a local build of the cited formalization, not a Comparator challenge-definition comparison.

The ordinary relative-consistency conclusion still uses the first-order interpretation in the manuscript's Section 8. Ambient Lean choice is part of the metatheory; it must not be confused with choice holding inside the modeled universe. Nor does a successful build of the original construction formally verify every new corollary in this report.

Two small additional developments were checked: \path{SixCycleAudit.lean} formalizes the finite cyclic-overlap and copied-assignment lemmas used in the BPI proof; \path{FiniteHeightAudit.lean} formalizes a finite-height combinatorial part of the order-forcing chain-condition argument. Their scope files identify the boundaries. Neither file formalizes the full forcing theorem, the entire six-cycle argument, the preservation theorem, the NDS deduction or the new exponentiation theorem.

\subsection{Mathematical and finite checks}
Separate AI-assisted reviews checked the finite-basis proof, exact comparison and lattice arguments, monoid-action and nonsplitting arguments, the six-cycle proof and its support/Cohen compatibility, parameter variants, cardinal-function coding, and the stated finite-order-extension results. The countable-member choice consequence received a second independent AI reconstruction during this consolidation. These checks can expose gaps and mismatched quantifiers, but they are not independent human peer review.

Finite computational checks validate the small graph gadget and the stated finite-order compatibility patterns; the saved files include the exact scopes and counts. Such finite checks do not replace the infinite forcing argument. Compilation, reference checks and visual PDF inspection concern the document, not mathematical truth.

\subsection{Novelty boundaries}
The literature record is a bounded search dated 6--7 October 2026. It includes a keyword search of the released text-source corpus across 722 OpenAI manuscripts (9,015 retrieved text-source files, hash-checked), followed by review of candidate matches. This was not a semantic proof-by-proof reading of all manuscripts. The target PP construction was examined directly.
\begin{longtable}{>{\raggedright\arraybackslash}p{.30\linewidth}>{\raggedright\arraybackslash}p{.62\linewidth}}
\toprule
Claim or subject & Attribution boundary\\\midrule\endhead
PP with failure of AC & The OpenAI construction; not a new model constructed by this report.\\
CSB* and WPP & Already explicit in the cited source, Remark 7.3.\\
Finite basis, NDS, finite cardinal antichains in $W$ & No earlier explicit application located in the bounded search. The normal-form mechanism is extracted from the source, and the ordinal-vector argument is classical. No priority certificate.\\
BPI/OEP failure in $W$ & No explicit statement located in the checked target source or bounded follow-up search. The new argument is a candidate additional application, pending specialist assessment.\\
Cardinal lattice and additive absorption & Related weaker consistency results are classical; broad novelty must not be inferred from a fresh proof or new notation.\\
Comparable power-set cardinals without AC & Already obtained in the first Cohen model in the cited Glazer discussion; not new as a bare consistency result.\\
Unequal exponents with equal powers at every base & Older unequal mutually surjective examples provide this phenomenon without PP. The present argument uses a different mechanism, with a well-orderable exponent and PP. Its exact prior status remains subject to further checking.\\
Real Hartogs/ordinal-choice package & Earlier comparison models exist. The new finite-order forcing analysis is the particular object studied here; no first-ever claim for its broad choice/real-cardinal package.\\
RPP versus BPI & Classical comparisons precede this report; merely obtaining the ordinary RPP here is not a new general implication.\\
\bottomrule
\end{longtable}

The distinction between injection and surjection, members and indices, set families and sequences, and $\ZF$ and $\ZFA$ is essential when searching for prior results. A statement under a different formulation may require a choice principle to transfer. The source-conflict records retain cases where a later introduction appeared to contradict older published results; no new theorem is based on the disputed historical claim. In particular, this report does not advertise consistency of $\mathsf{DC}+\mathsf{BPI}+\neg\AC$ as newly obtained.

\subsection{Open questions retained at the cut-off}
\begin{enumerate}
\item Does $\PP$ imply $\NDS$ or $\FB$ in arbitrary $\ZF$ models? Does NDS or FB imply PP? The model-specific positive package settles none of these unrestricted directions.
\item Is there a different pure model of $\PP+\mathsf{BPI}+\neg\AC$? Does $\SVC(A)+Q(A)+\ACWO+\mathsf{BPI}$ force choice? The negative BPI theorem for $W$ alone does not answer either question.
\item In $N=W[G]$, does $\mathbb R\inj A$ hold? Are all non-well-orderable images of $A$ equinumerous with $A$? Do PP, NDS, FB or BPI hold? Preserve the distinction between the old-set comparisons already proved and these global new-set questions.
\item Can the new real-sized bound on directed availability ideals replace the original ground-cardinal bound in the coherent quotient argument? A coding bound alone does not give the needed supported matching.
\item What are the exact two-colour triple Ramsey thresholds? The finite-colour spectrum and the failure of fixed coarsenings to supply the desired binary obstruction do not settle this.
\item Can an unordered ordinary-FM generator realize the exact coset criterion, RPC plus the missing dichotomy, or the finite-prototype criterion? The many exclusions are architecture-specific and do not establish a universal impossibility theorem.
\item Which exact all-base exponent-cancellation or exponent-separation principles are equivalent to AC or to known weak-choice principles? The present model refutes the stated principles, and the recorded implications are limited to those proved above.
\end{enumerate}

\subsection{Files, authorship and publication record}
The public package supplies this PDF, its complete LaTeX source, supporting mathematical notes, the two auxiliary Lean developments with check records, and a machine-readable source map with checksums. The evidence is grouped by original-model results, earlier FM investigations, source verification and novelty records. Chronological notes may retain an earlier unresolved status that was later settled; this report is the status guide, and the filenames cited in each chapter lead to the detailed arguments.

The original imported PDF and TeX remain unchanged in the local workspace. They are not included in the public evidence archive, which contains the mathematical follow-up records rather than the raw imported conversation or account/usage data. OpenAI's third-party manuscript and Lean source distribution are linked by pinned revision, not relicensed or republished here.

The earlier focused finite-basis paper remains available; this report extends the scope rather than overwriting its historical public version. The Weak Choice Atlas website is a separate software/catalogue output. Its search, comparison and provenance checks do not certify every mathematical result, and publication of this report does not imply that all its new records have already been added to that website.

This version is a preliminary research release by Frederico Can\c{c}ado, with substantial Codex assistance disclosed. A public GitHub commit documents this version and date; it does not certify correctness or priority. No DOI, journal acceptance, or independent human verification is asserted. Original expression is released under the repository's CC BY 4.0 terms to the extent the contributor holds rights in it. Unproved questions do not reserve their eventual solutions.

\clearpage\begin{thebibliography}{99}\small\raggedright
\bibitem{BM} B. Banaschewski and G. H. Moore, \emph{The dual Cantor--Bernstein theorem and the partition principle}, Notre Dame Journal of Formal Logic \textbf{31} (1990), no.~3, 375--381. \href{https://doi.org/10.1305/ndjfl/1093635502}{doi:10.1305/ndjfl/1093635502}.

\bibitem{Blass} A. Blass, \emph{Injectivity, projectivity, and the axiom of choice}, Transactions of the American Mathematical Society \textbf{255} (1979), 31--59. \href{https://doi.org/10.1090/S0002-9947-1979-0542870-6}{doi:10.1090/S0002-9947-1979-0542870-6}.

\bibitem{BK} A. Blass and D. Kulshreshtha, \emph{Cardinal Well-foundedness and Choice}, arXiv:2310.09643v3, 21 August 2026. \url{https://arxiv.org/abs/2310.09643v3}.

\bibitem{daSilva} S. G. da Silva, \emph{The Axiom of Choice and the Partition Principle from Dialectica Categories}, Logic Journal of the IGPL \textbf{29} (2021), no.~5, 783--797. \href{https://doi.org/10.1093/jigpal/jzaa023}{doi:10.1093/jigpal/jzaa023}.

\bibitem{FO} D. Feldman and M. Orhon, with an appendix by A. Blass, \emph{Generalizing Hartogs' Trichotomy Theorem}, arXiv:0804.0673 (2008). \url{https://arxiv.org/abs/0804.0673}. See Theorem 1 and the discussion of infinite trichotomy.

\bibitem{HT} P. Howard and E. Tachtsis, \emph{No decreasing sequence of cardinals}, Archive for Mathematical Logic \textbf{55} (2016), 415--429. \href{https://doi.org/10.1007/s00153-015-0472-5}{doi:10.1007/s00153-015-0472-5}.

\bibitem{Karagila} A. Karagila, \emph{Embedding orders into the cardinals with $\mathsf{DC}_\kappa$}, Fundamenta Mathematicae \textbf{226} (2014), no.~2, 143--156. \href{https://doi.org/10.4064/fm226-2-4}{doi:10.4064/fm226-2-4}.

\bibitem{OpenAI} OpenAI, \emph{The Partition Principle does not imply Choice}, manuscript dated 24 September 2026, repository revision checked 7 October 2026. Commit \texttt{adc7f1241b42e322a6451854ab7e4b4c146bf78a}. \href{https://github.com/openai/math/tree/adc7f1241b42e322a6451854ab7e4b4c146bf78a/preprints/The-Partition-Principle-does-not-imply-Choice-September-24-2026}{Pinned manuscript and source}.

\bibitem{Pelc} A. Pelc, \emph{On Some Weak Forms of the Axiom of Choice in Set Theory}, Bulletin de l\'Acad\'emie Polonaise des Sciences, S\'erie des Sciences Math\'ematiques, Astronomiques et Physiques \textbf{26} (1978), 585--589. The results used for historical context are discussed in \cite[Section 10]{BK}.

\bibitem{RyanSmith} C. Ryan-Smith, \emph{Local reflections of choice}, Acta Mathematica Hungarica \textbf{176} (2025), no.~1, 244--257. \href{https://doi.org/10.1007/s10474-025-01533-3}{doi:10.1007/s10474-025-01533-3}. Author manuscript: \url{https://academic.calliope.mx/files/papers/2412.13785.pdf}. Section references in this note refer to the checked author manuscript.

\bibitem{Tachtsis2018} E. Tachtsis, \emph{On the comparability of cardinals in the absence of the axiom of choice}, Fundamenta Mathematicae \textbf{242} (2018), no.~3, 247--266. \href{https://doi.org/10.4064/fm376-11-2017}{doi:10.4064/fm376-11-2017}.

\bibitem{Tachtsis2024} E. Tachtsis, \emph{No Decreasing Sequence of Cardinals in the Hierarchy of Choice Principles}, Notre Dame Journal of Formal Logic \textbf{65} (2024), no.~3, 311--331. \href{https://doi.org/10.1215/00294527-2024-0020}{doi:10.1215/00294527-2024-0020}.

\bibitem{FelgnerTruss} U. Felgner and J. K. Truss, \emph{The independence of the Prime Ideal Theorem from the Order-Extension Principle}, J. Symbolic Logic \textbf{64} (1999), 199--215. \href{https://doi.org/10.2307/2586759}{doi:10.2307/2586759}.

\bibitem{Higman} G. Higman, \emph{Ordering by Divisibility in Abstract Algebras}, Proc. London Math. Soc. (3) \textbf{2} (1952), 326--336. \href{https://doi.org/10.1112/plms/s3-2.1.326}{doi:10.1112/plms/s3-2.1.326}.

\bibitem{Jech} T. Jech, \emph{The Axiom of Choice}, North-Holland, 1973. Example 2.3.2 and Section 7.2. \url{https://gwern.net/doc/math/1973-jech-theaxiomofchoice.pdf}.

\bibitem{Rado} R. Rado, \emph{Partial well-ordering of sets of vectors}, Mathematika \textbf{1} (1954), 89--95. \href{https://doi.org/10.1112/S0025579300000565}{doi:10.1112/S0025579300000565}.

\bibitem{Glazer} E. Glazer, answer to D. van der Zypen, \emph{Comparability of power sets and (AC)}, MathOverflow, 30 September 2025. \url{https://mathoverflow.net/a/501058}. Public mathematical communication, not a refereed article.

\bibitem{BSWY} M. Beriashvili, R. Schindler, L. Wu and L. Yu, \emph{Hamel bases and well-ordering the continuum}, Proc. Amer. Math. Soc. 146 (2018), 3565--3573. \url{https://doi.org/10.1090/proc/14010}.

\bibitem{ICFquestion} A. Karagila, \emph{Does $2^X=2^Y\Rightarrow X=Y$ imply the Axiom of Choice?}, MathOverflow, 17 November 2011, with subsequent answers and discussion. \url{https://mathoverflow.net/questions/81204}. Historical question and terminology; current open status is not inferred from the date of the page.

\bibitem{Hamkins} J. D. Hamkins, answer to A. Karagila, \emph{Does $2^X=2^Y\Rightarrow X=Y$ imply the Axiom of Choice?}, MathOverflow, 24 August 2012. \url{https://mathoverflow.net/a/105424}.

\bibitem{LT} A. Lindenbaum and A. Tarski, \emph{Communication sur les recherches de la th\'eorie des ensembles}, Comptes Rendus des S\'eances de la Soci\'et\'e des Sciences et des Lettres de Varsovie, Classe III, 19 (1926), 299--330, No.~82(A5), p.~312. Original scan: \url{https://rcin.org.pl/Content/50065}.

\bibitem{Grochow} J. A. Grochow, \emph{Polynomial-Time Axioms of Choice and Polynomial-Time Cardinality}, Theory of Computing Systems 67 (2023), 627--669. \url{https://doi.org/10.1007/s00224-023-10118-y}. Appendix A.3 reproduces the related Rubin--Rubin strict-reflection formulation.

\bibitem{KS} A. Karagila and J. Schilhan, \emph{Intermediate models and Kinna--Wagner Principles}, Proc. Amer. Math. Soc. 154 (2026), 393--403. \url{https://doi.org/10.1090/proc/17425}.

\bibitem{Blass1977} A. Blass, \emph{Ramsey's theorem in the hierarchy of choice principles}, Journal of Symbolic Logic \textbf{42} (1977), no.~3, 387--390. \href{https://doi.org/10.2307/2272866}{doi:10.2307/2272866}.

\bibitem{ErdosRado} P. Erd\H{o}s and R. Rado, \emph{A partition calculus in set theory}, Bulletin of the American Mathematical Society \textbf{62} (1956), 427--489. Theorem~39, pp.~467--469, especially clause~(iii). Original paper: \url{https://www.renyi.hu/~p_erdos/1956-02.pdf}.

\bibitem{BHT} J. E. Baumgartner, A. Hajnal and S. Todor\v{c}evi\'c, \emph{Extensions of the Erd\H{o}s--Rado theorem}, manuscript, arXiv:math/9311207 (1993). Theorem~2.1, manuscript p.~4. \url{https://arxiv.org/abs/math/9311207}.
\end{thebibliography}

\end{document}
